% Paper 1 --- main file.
% Build (from this directory): pdflatex -output-directory=.build main; in .build: BIBINPUTS="..;." bibtex main;
% then pdflatex -output-directory=.build main twice more. .build/ is ignored by git.
\documentclass[aps,prb,reprint,amsmath,amssymb,superscriptaddress,floatfix,nofootinbib]{revtex4-2}

\usepackage[T1]{fontenc}
\usepackage[utf8]{inputenc}
\usepackage{lmodern}
\usepackage{amsthm}
\usepackage{bm}
\usepackage{mathtools}
\usepackage{xcolor}
\usepackage[colorlinks=true,linkcolor=blue!50!black,citecolor=blue!50!black,urlcolor=blue!50!black]{hyperref}

% ------------------------------------------------------------------ markers and labels
% Items whose content is not yet established print in red.
\newcommand{\PENDING}[1]{{\color{red}\textbf{[PENDING:}~#1\textbf{]}}}
% Status labels of statements (proved / derived under assumptions / open / numerically demonstrated).
\newcommand{\status}[1]{\textup{\textsc{#1}}}

\hyphenation{thermo-re-flec-tance semi-def-i-nite Krum-hansl Cat-ta-neo ad-mis-si-ble}

% ------------------------------------------------------------------ statements
\theoremstyle{plain}
\newtheorem{theorem}{Theorem}
\newtheorem{proposition}{Proposition}
\newtheorem{lemma}{Lemma}
\newtheorem{corollary}{Corollary}
\theoremstyle{definition}
\newtheorem{definition}{Definition}
\newtheorem{example}{Example}
\theoremstyle{remark}
\newtheorem{remark}{Remark}

% ------------------------------------------------------------------ notation
\newcommand{\tr}{^{\mathsf T}}
\newcommand{\Cop}{\mathcal C}
\newcommand{\Hsp}{\mathcal H}
\newcommand{\ee}{\mathbf e}
\newcommand{\jj}{\mathbf j}
\newcommand{\xx}{\mathbf x}
\newcommand{\yy}{\mathbf y}
\newcommand{\tmem}{\tau_{\rm mem}}
\newcommand{\tstat}{\tau_{\rm stat}}

\begin{document}

\title{What static transport data determine about a conserving phonon closure}

\author{Koussay Mansouri}
\affiliation{Centre de Recherches et des Technologies de l'\'Energie (CRTEn), Technop\^ole de Borj C\'edria, Tunisia}
\affiliation{Centre National des Sciences et Technologies Nucl\'eaires (CNSTN), P\^ole Technologique de Sidi Thabet, Tunisia}
\affiliation{INSA Hauts-de-France, Universit\'e Polytechnique Hauts-de-France, Valenciennes, France}

\date{\today}

\begin{abstract}
Non-Fourier interpretations of thermoreflectance experiments on thin films attach kinetic meaning to
the relaxation time and the nonlocal length of hydrodynamic heat-flux laws. We ask what static
transport data---mode lifetimes, collision invariants and the dc thermal conductivity---determine
about the dynamic part of an energy-conserving closure of the linearized phonon Boltzmann equation,
and which further operator or measurement information is required. In entropy coordinates the
Cattaneo time of the closure is the first memory moment $\tmem=-K'(0)/K(0)$ of the current response.
Static data bound $\tmem$ from below by $K(0)/|\jj|^2$, but in the class of positive-semidefinite
collision operators they do not bound it from above: an explicit family keeps all lifetimes, the
invariant and $K(0)$ fixed while $\tmem\to\infty$. If the collision events are fixed by the harmonic
spectrum and only their nonnegative rates are unknown, a convex-hull theorem of Ben-Tal and Teboulle
gives $\tmem\le M^2K(0)$ with a sharp, rate-independent constant $M$, and a per-event lower bound on
the rates bounds $\tmem$ without any data. In the relaxation-time approximation the memory moment
diverges in the continuum when the infrared exponent $\alpha$ of the acoustic rates satisfies
$2\alpha\ge d$, while the conductivity stays finite for $\alpha<d$; for event operators the same
divergence follows under explicit infrared hypotheses, and then $M^2K(0)$ diverges under
Brillouin-zone refinement. For a one-dimensional front-face measurement on a Cattaneo or
Guyer--Krumhansl film on a substrate, the depth-scaling gauge of the Fourier problem persists: the film
enters only through $(\xi_1,b,\tau_R,\tau_\ell)$, so neither the nonlocal length nor the closure
coefficient is separately identifiable. \PENDING{abstract sentences on the AlN film-regularized
memory-time ranges, FDTR detectability and the rank theorem}
\end{abstract}

\maketitle

% Section I — Introduction
\section{Introduction}\label{sec:intro}

Thermoreflectance in the time domain (TDTR) and in the frequency domain (FDTR) is a standard probe of
heat transport in thin films and across interfaces~\cite{Cahill2004,Schmidt2009}. When the heated volume or the
modulation period approaches phonon mean free paths or lifetimes, the data depart from the predictions
of Fourier's law: apparent conductivities depend on the modulation frequency~\cite{KohCahill2007,Regner2013}
and on the size and anisotropy of the heating~\cite{WilsonCahill2014}, and such departures have been
inverted for mean-free-path or relaxation-time spectra~\cite{Minnich2012,Forghani2016,Chen2025}. A second line
of analysis replaces Fourier's law by a constitutive law with a heat-flux relaxation
time~$\tau_R$---the Cattaneo law~\cite{JosephPreziosi1989}---or, in addition, a nonlocal
length~$\ell$---the Guyer--Krumhansl (GK) equation~\cite{GuyerKrumhansl1966a,GuyerKrumhansl1966b,GuoWang2015}---and
interprets thermoreflectance signals through these
parameters~\cite{Beardo2020,HennessyMyers2021,CamachoDeLaRosa2025}. Their values are then read
kinetically: $\tau_R$ as a resistive (momentum-relaxing) time and $\ell$ as a length set by normal
processes.

Aluminum nitride is a natural test case for this inference. Its thermal conductivity is high---about
$321~\mathrm{W\,m^{-1}K^{-1}}$ at room temperature was measured by TDTR on films grown to $18$ and
$22.5~\mu\mathrm{m}$~\cite{Cheng2020}, a value that cannot be transferred to a submicrometer film without
a measurement---and first-principles transport workflows for it have been benchmarked against one
another~\cite{McGaughey2025}; Callaway-type descriptions of its phonon transport have been examined
against first-principles rates~\cite{MaLiLuo2014}. A first-principles route from AlN to a
thin-film thermoreflectance prediction has three steps: (i)~static transport data---mode lifetimes
and the bulk conductivity tensor---are computed from harmonic and anharmonic force constants;
(ii)~an energy-conserving closure of the linearized Boltzmann equation is constructed from them;
(iii)~the closure enters a model of the film response. The closure coefficients that govern the
dynamics, above all the Cattaneo time, are not static quantities. This paper addresses the question:
\emph{what do static transport data determine about a conserving dynamic closure, and what further
operator or measurement information is needed?}

Several ingredients are established and are used here as attributed background. For the grey Callaway
model~\cite{Callaway1959}, the conserving closure has grad-div coefficient $1/3$ instead of the
coefficient $2$ of the original GK equation; Guo and Wang give it, including the combined time
$\tau_c$, in Eq.~(7) of Ref.~\cite{GuoWang2022}, Sendra \emph{et~al.} derive and discuss it~\cite{Sendra2021,Sendra2022},
and the coefficient $2$ has been attributed to an approximate inverse of the normal-process collision
matrix~\cite{HardyAlbers1974,GuoWang2022}; see also Ref.~\cite{Struchtrup2025}. The eigenvectors of the symmetrized collision operator
(relaxons) express the conductivity as a sum over relaxation modes~\cite{CepellottiMarzari2016}, and
conductivity and thermal viscosity are carried by relaxons of opposite parity~\cite{Simoncelli2020}.
Reductions of collision dynamics to a few collective variables are exact only under lumpability
conditions~\cite{WeiKuo1969,Aoki1968,LiRabitz1989,Buchholz1994} and otherwise generate memory terms of
projection type~\cite{Zwanzig1961,Mori1965}, on which a recent statistical theory of heat conduction with
spatiotemporal memory kernels is built~\cite{ZengDong2026}; descriptions built on diagonal rates alone can fail for
hydrodynamic transport~\cite{MalviyaRavichandran2023}. For Fourier films, Krapez and Rigollet showed that
thickness, diffusivity and conductivity of a coating are not separately identifiable from front-face
data, whatever the frequencies and noise level~\cite{KrapezRigollet2017}.

Within this background the paper makes the following contributions. Each is stated with its
hypotheses and one of the status labels defined below.
\begin{enumerate}
\item In entropy coordinates the Cattaneo time of a conserving closure, defined by matching a single-pole
law to the dynamic conductivity at zero frequency, is the first memory moment
$\tmem=-K'(0)/K(0)=|\Cop^+\jj|^2/\jj\tr\Cop^+\jj$ of the current response (Sec.~\ref{sec:kinetics});
in the grey Callaway model it equals $\tau_R$ exactly. This is a definition and an elementary check, not a new result.
\item Static data (lifetimes, invariants and the dc response) bound $\tmem$ from below by
$\tstat=K(0)/|\jj|^2$ (Proposition~\ref{prop:lower}) but, in the class of positive-semidefinite
operators, not from above (Theorem~\ref{thm:psd}). The mathematics is elementary: it is an instance of the
indeterminacy of a Stieltjes function from finitely many of its moments~\cite{Akhiezer2020,Schmudgen2017}.
We have not found an explicit identifiability statement of this kind for phonon kinetics in the
literature searched.
\item If the events are fixed by the harmonic spectrum and only their nonnegative rates are unknown,
$\tmem\le M^2K(0)$ with a sharp, rate-independent constant~$M$ (Theorem~\ref{thm:cone} and
Proposition~\ref{prop:sharp}). The inequality is a known consequence of the convex-hull theorem of Ben-Tal and
Teboulle for weighted least squares~\cite{BenTalTeboulle1990}; its application to the memory of kinetic
operators is potentially novel within the literature searched. A per-event lower bound on the rates bounds
$\tmem$ without any data (Proposition~\ref{prop:prior}).
\item In the relaxation-time approximation (RTA) the memory moment diverges in the continuum when the
infrared exponent of the acoustic rates satisfies $2\alpha\ge d$, while the conductivity is finite for
$\alpha<d$; for event operators the divergence follows under explicit hypotheses on the infrared
rates and solution (Theorem~\ref{thm:ir}), and then $M^2K(0)$ diverges under mesh refinement. No novelty is claimed
for this criterion; an analogous coexistence of a finite conductivity with divergent higher moments, of
the free-path distribution, underlies the superdiffusive description of quasiballistic transport in
alloys~\cite{Vermeersch2015a,Vermeersch2015b}.
\item For a one-dimensional front-face measurement on a Cattaneo or GK film on a substrate, the
depth-scaling gauge of the Fourier problem persists: the film enters only through
$(\xi_1,b,\tau_R,\tau_\ell)$, so neither $\ell$ nor the closure coefficient is separately identifiable
(Proposition~\ref{prop:gauge} and Corollary~\ref{cor:beta}). The Fourier gauge is
known~\cite{KrapezRigollet2017}; its extension to non-Fourier laws is elementary dimensional analysis and is
potentially novel within the literature searched. Whether the four combinations are themselves identifiable is
the subject of a theorem stated in Sec.~\ref{sec:frontface}: \PENDING{Theorem: generic rank four and degenerate set}.
\item \PENDING{AlN first-principles results: film-regularized memory-time ranges (Sec.~\ref{sec:aln})}
\PENDING{FDTR detectability of those ranges (Sec.~\ref{sec:fdtr})}
\end{enumerate}

The paper does not claim a value of the bulk memory time of AlN, an experimental validation, or any
part of the conserving closure itself. The Fourier resonance $\tau_R=\tau_\ell$ used in
Sec.~\ref{sec:frontface} was named by Both \emph{et~al.}~\cite{Both2016} and analyzed by
Kov\'acs~\cite{Kovacs2018}.

\emph{Status labels.} \status{proved}: a complete proof is given. \status{derived under assumptions}:
the derivation is complete given hypotheses that are stated but not proved (for instance a property of a
specific operator, or a model of the experiment). \status{open}: neither proved nor refuted. Numerical
results in Secs.~\ref{sec:aln} and \ref{sec:fdtr} will carry the label
\status{numerically demonstrated} together with their error estimates. Statements of novelty are
restricted to the literature we searched in October 2026: Crossref and OpenAlex (about fifty queries
each), arXiv full text, publisher and repository pages, and the forward citations of
Refs.~\cite{BenTalTeboulle1990,Stewart1989}.

Section~\ref{sec:kinetics} sets up the kinetic framework, the memory moment and the two operator
classes. Section~\ref{sec:identifiability} states and proves what static data do and do not determine.
Section~\ref{sec:frontface} formulates the front-face measurement model and its gauge.
Sections~\ref{sec:aln}--\ref{sec:conclusions} are reserved for the AlN results, the detectability
analysis and the conclusions. The notation is collected in Appendix~\ref{app:notation}.

% Section II — Linearized kinetics in entropy coordinates
\section{Linearized kinetics in entropy coordinates}\label{sec:kinetics}

\subsection{Modes, entropy coordinates and the collision operator}\label{sec:modes}

We work with a finite, discretized description of the phonon gas; continuum limits are taken explicitly
in Sec.~\ref{sec:ir}.
\begin{enumerate}
\item[(K1)] A crystal at uniform temperature $T>0$ is described on a finite $\Gamma$-centered mesh of
$N_{\mathbf k}$ wave vectors, invariant under $\mathbf k\mapsto-\mathbf k$; the crystal volume is
$\mathcal V=N_{\mathbf k}V_{\rm cell}$. The modes $\nu=(\mathbf k,s)$, $\nu=1,\dots,n$, have energies
$\varepsilon_\nu=\hbar\omega_\nu>0$ (zero-frequency modes at $\Gamma$ are excluded), group velocities
$\mathbf v_\nu$ and Bose occupations $f^0_\nu=[\exp(\varepsilon_\nu/k_BT)-1]^{-1}$.
\item[(K2)] The occupations obey the linearized Boltzmann equation
$\partial_t\delta f_\nu+\mathbf v_\nu\cdot\nabla\delta f_\nu=-(\Omega\,\delta f)_\nu$, and the collision
operator satisfies detailed balance and the $H$ theorem: with $\mathsf W=\mathrm{diag}(w_\nu)$ and
$w_\nu=[f^0_\nu(1+f^0_\nu)]^{1/2}$, the matrix $\Omega\mathsf W^2$ is symmetric and positive semidefinite.
\end{enumerate}
In the entropy coordinates
\begin{equation}
y_\nu=\frac{\delta f_\nu}{\mathcal V^{1/2}w_\nu}
\label{eq:entropycoords}
\end{equation}
the kinetic equation reads
\begin{equation}
\partial_t\yy+\mathsf V\cdot\nabla\yy=-\Cop\,\yy,\qquad \Cop=\mathsf W^{-1}\Omega\,\mathsf W=\Cop\tr\succeq0,
\label{eq:kinetic}
\end{equation}
with $\mathsf V=\mathrm{diag}(\mathbf v_\nu)$. Full linearized collision operators of this kind are
computed from first principles~\cite{Chaput2013,Fugallo2013,CepellottiMarzari2016}. The name of the
coordinates refers to the entropy density: expanding the
Bose entropy $s_B(f)=k_B[(1+f)\ln(1+f)-f\ln f]$ of each mode to second order gives
\begin{equation}
\delta s=\frac{\delta E}{T}-\frac{k_B}{2}\,|\yy|^2+O(|\yy|^3),\qquad \delta E=\ee\tr\yy,
\label{eq:entropy}
\end{equation}
where $e_\nu=\mathcal V^{-1/2}w_\nu\varepsilon_\nu$, because $s_B'(f^0_\nu)=\varepsilon_\nu/T$ and
$s_B''(f^0_\nu)=-k_B/w_\nu^2$. At fixed energy the Euclidean norm of $\yy$ measures the entropy deficit,
and collisions produce entropy at the rate $k_B\,\yy\tr\Cop\yy\ge0$. Energy conservation in collisions
gives $\Cop\ee=0$.
\begin{enumerate}
\item[(K3)] The energy is the only collision invariant: $\ker\Cop=\mathrm{span}\{\ee\}$.
\end{enumerate}
Hypothesis (K3) excludes exact conservation of crystal momentum (umklapp and other resistive processes
are present). If further invariants exist, everything below holds with $\Hsp$ replaced by the
orthogonal complement of all invariants. We write $\Hsp=\ee^\perp$, $n_{\Hsp}=\dim\Hsp=n-1$ and
$\Pi_{\Hsp}$ for the orthogonal projector onto $\Hsp$; $\Cop$ maps $\Hsp$ into itself and is positive
definite there.

The heat-flux density along the Cartesian direction $i$ is $q_i=\jj^{(i)\mathsf T}\yy$ (a flux per unit
area, in $\mathrm{W\,m^{-2}}$), with the current vector
\begin{equation}
j^{(i)}_\nu=\mathcal V^{-1/2}\,w_\nu\,\varepsilon_\nu\,v_{\nu,i}.
\label{eq:current}
\end{equation}
Time reversal $\nu=(\mathbf k,s)\mapsto\bar\nu=(-\mathbf k,s)$ leaves $\ee$ invariant and reverses
$\jj^{(i)}$, so $\jj^{(i)}\in\Hsp$. The harmonic quantities
\begin{equation}
\rho c=\frac{|\ee|^2}{k_BT^2},\qquad
\Lambda_i=\frac{|\jj^{(i)}|^2}{k_BT^2}=\frac1{\mathcal V}\sum_\nu c_\nu v_{\nu,i}^2,
\label{eq:harmonic}
\end{equation}
are the volumetric heat capacity and a ballistic weight; $c_\nu=\varepsilon_\nu^2w_\nu^2/(k_BT^2)$ is the
heat capacity of mode $\nu$. The first identity follows from the response $\yy=\ee\,\delta T/(k_BT^2)$
of the equilibrium to a temperature change $\delta T$.

\subsection{Dynamic conductivity and memory kernel}\label{sec:response}

Consider the response to a spatially uniform temperature gradient $\nabla T(t)$ switched on at $t=0$
(the zero-wave-vector limit). Write the occupations as the local equilibrium $f^0_\nu(T(\mathbf r,t))$,
with the local temperature fixed by the energy density, plus a deviation whose entropy coordinates
$\yy$ lie in $\Hsp$. The streaming of the local equilibrium,
$\mathbf v_\nu\cdot\nabla f^0_\nu=w_\nu^2\varepsilon_\nu\,\mathbf v_\nu\cdot\nabla T/(k_BT^2)$, acts as a
source with entropy coordinates $\sum_k\jj^{(k)}\partial_kT/(k_BT^2)\in\Hsp$; the time derivative of the
local equilibrium is proportional to $\ee$ and drops out on projection onto $\Hsp$ (its projection onto
$\ee$ is the energy balance). For a spatially uniform deviation,
\begin{equation}
\partial_t\yy=-\Cop\yy-\frac{1}{k_BT^2}\sum_k\jj^{(k)}\,\partial_kT(t),\qquad \yy(0)=0.
\label{eq:drive}
\end{equation}
Its solution gives the heat flux in the form of a constitutive law with memory~\cite{GurtinPipkin1968},
\begin{equation}
q_i(t)=-\int_0^t\kappa_{ik}(t-t')\,\partial_kT(t')\,dt',\qquad
\kappa_{ik}(t)=\frac{\Phi_{ik}(t)}{k_BT^2},
\label{eq:memorylaw}
\end{equation}
with the current autocorrelation function $\Phi_{ik}(t)=\jj^{(i)\mathsf T}e^{-\Cop t}\jj^{(k)}$. Under the
time convention $e^{+i\omega t}$ of this paper, the Laplace variable is $p$ and the harmonic regime is
$p=i\omega$. The transform of Eq.~\eqref{eq:memorylaw} defines the dynamic conductivity
\begin{equation}
\lambda_{ik}(p)=\frac{K_{ik}(p)}{k_BT^2},\qquad K_{ik}(p)=\jj^{(i)\mathsf T}(p+\Cop)^{-1}\jj^{(k)}.
\label{eq:dynamicK}
\end{equation}
Because the currents lie in $\Hsp$, where $\Cop$ is invertible, $K_{ik}$ is analytic at $p=0$ and the
dc conductivity is $\lambda_{ik}=K_{ik}(0)/(k_BT^2)$ with $K_{ik}(0)=\jj^{(i)\mathsf T}\Cop^+\jj^{(k)}$;
$\Cop^+$ is the Moore--Penrose pseudoinverse. A current with a component along an invariant would add a
term proportional to $1/p$, and the pseudoinverse must not be used to hide it. From now on we fix one
direction and drop the index $i$: $\jj$, $K(p)$, $\Phi(t)$, $\lambda(p)$ and $\lambda=\lambda(0)$.

With the eigenpairs $(\Gamma_k,\mathbf u_k)$ of $\Cop$ restricted to $\Hsp$ ($\Gamma_k>0$; the relaxons
of Ref.~\cite{CepellottiMarzari2016}),
\begin{equation}
K(p)=\sum_{k=1}^{n-1}\frac{(\mathbf u_k\tr\jj)^2}{p+\Gamma_k}.
\label{eq:stieltjes}
\end{equation}
$K$ is the Stieltjes transform of the spectral measure of $\Cop|_{\Hsp}$ in the state $\jj$; its total
mass is $|\jj|^2=\Phi(0)$.

\subsection{The memory moment}\label{sec:memory}

\begin{definition}[Moments]\label{def:moments}
Under (K1)--(K3) and for $\jj\in\Hsp\setminus\{0\}$, let $\xx=\Cop^+\jj$ (the unique solution in $\Hsp$
of $\Cop\xx=\jj$), and define
\begin{equation}
\begin{gathered}
K_0=K(0)=\jj\tr\xx,\qquad N=|\xx|^2,\\
\tmem=-\frac{K'(0)}{K(0)},\qquad \tstat=\frac{K_0}{|\jj|^2}.
\end{gathered}
\label{eq:moments}
\end{equation}
\end{definition}

\begin{lemma}[Equivalent forms of the memory moment; \status{proved}]\label{lem:forms}
Under the hypotheses of Definition~\ref{def:moments},
\begin{equation}
\tmem=\frac{N}{K_0}=\frac{|\Cop^+\jj|^2}{\jj\tr\Cop^+\jj}
=\frac{\int_0^\infty t\,\Phi(t)\,dt}{\int_0^\infty\Phi(t)\,dt}
=\frac{\sum_k\lambda_k\Gamma_k^{-1}}{\sum_k\lambda_k},
\label{eq:forms}
\end{equation}
where $\lambda_k=(\mathbf u_k\tr\jj)^2/(k_BT^2\,\Gamma_k)$ is the contribution of relaxon $k$ to
$\lambda$. In particular $\Gamma_{\max}^{-1}\le\tmem\le\Gamma_{\min}^{-1}$, where $\Gamma_{\min}$ and
$\Gamma_{\max}$ are the extreme eigenvalues of $\Cop|_{\Hsp}$.
\end{lemma}

\begin{proof}
On $\Hsp$, $(p+\Cop)^{-1}$ is analytic at $p=0$ with derivative $-(p+\Cop)^{-2}$; hence
$-K'(0)=\jj\tr(\Cop^+)^2\jj=|\Cop^+\jj|^2$ because $\Cop^+$ is symmetric. In time,
$\int_0^\infty e^{-\Cop t}dt=\Cop^{+}$ and $\int_0^\infty t\,e^{-\Cop t}dt=(\Cop^{+})^2$ on $\Hsp$. The
last form follows from Eq.~\eqref{eq:stieltjes}: $K_0=\sum_k(\mathbf u_k\tr\jj)^2\Gamma_k^{-1}$ and
$N=\sum_k(\mathbf u_k\tr\jj)^2\Gamma_k^{-2}$. A weighted mean of the $\Gamma_k^{-1}$ lies between their
extremes.
\end{proof}

Thus $\tmem$ is the first moment of the normalized memory kernel $\kappa(t)/\lambda$ and the
conductivity-weighted mean relaxon lifetime. In the RTA, where $\Cop$ is replaced by
$\mathrm{diag}(r_\nu)$ (an operator that does not conserve energy),
$\tmem^{\rm RTA}=\sum_\nu j_\nu^2r_\nu^{-2}/\sum_\nu j_\nu^2r_\nu^{-1}$ and
$\tstat^{\rm RTA}=\sum_\nu j_\nu^2r_\nu^{-1}/\sum_\nu j_\nu^2$.

\subsection{The Cattaneo time of a conserving closure}\label{sec:cattaneo}

For spatially uniform fields, Cattaneo's law $\tau_R\partial_tq+q=-\lambda\,\partial_iT$ is
Eq.~\eqref{eq:memorylaw} with the exponential kernel $\kappa_C(t)=(\lambda/\tau_R)e^{-t/\tau_R}$, that
is, with the single pole $\lambda_C(p)=\lambda/(1+p\tau_R)$~\cite{JosephPreziosi1989,GurtinPipkin1968}.

\begin{definition}[Cattaneo time]\label{def:cattaneo}
The Cattaneo time of a conserving operator $\Cop$ (for a given direction) is the relaxation time of the
single pole that matches $\lambda(p)$ and $\lambda'(p)$ at $p=0$. Since $\lambda_C(0)=\lambda$ and
$\lambda_C'(0)=-\lambda\tau_R$, it equals $\tmem$.
\end{definition}

The match is exact for all $p$ if and only if $\Phi$ is a single exponential, that is, if and only if
$\jj$ is an eigenvector of $\Cop$. The matched pole reproduces the slope of $\lambda(p)$ at $p=0$, not
the response in a finite frequency band; Theorem~\ref{thm:psd}(v) below shows that the two can be
unrelated.

The grey conserving closure illustrates Definition~\ref{def:cattaneo}. In the Callaway
model~\cite{Callaway1959} with a common linear isotropic dispersion $\varepsilon=\hbar v|\mathbf k|$
and constant times $\tau_N,\tau_R>0$, the linearized operator in entropy coordinates is
\begin{equation}
\Cop_{\rm Cal}=\tau_N^{-1}(I-\Pi_{\rm N})+\tau_R^{-1}\,\Pi_{\Hsp},
\label{eq:callaway}
\end{equation}
where $\Pi_{\rm N}$ is the orthogonal projector onto the span of $\ee$ and of the three vectors
$\mathbf P_x$, $\mathbf P_y$, $\mathbf P_z$, with $P_{i,\nu}=\mathcal V^{-1/2}w_\nu\hbar k_{\nu,i}$,
which represent crystal momentum: the displaced
equilibrium that matches the energy and momentum of a state $\yy$ is its orthogonal projection onto that
span. We take the mesh symmetric under the cubic point group, so that $\ee,\mathbf P_x,\mathbf P_y,\mathbf P_z$ are
mutually orthogonal.

\begin{proposition}[Grey conserving closure; \status{derived under assumptions}]\label{prop:callaway}
Under the assumptions just stated, $\jj^{(i)}=v^2\mathbf P_i$ and $\Cop_{\rm Cal}\jj^{(i)}=\jj^{(i)}/\tau_R$.
Hence $\lambda(p)=\lambda/(1+p\tau_R)$ exactly, with $\lambda=\rho c\,v^2\tau_R/3$, and
$\tmem=\tstat=\tau_R$. The mode relaxation rates are
$(\Cop_{\rm Cal})_{\nu\nu}=\tau_N^{-1}+\tau_R^{-1}-O(n^{-1})$.
\end{proposition}

\begin{proof}
For $\varepsilon=\hbar v|\mathbf k|$, $\varepsilon_\nu v_{\nu,i}=\hbar v^2k_{\nu,i}$, so
$\jj^{(i)}=v^2\mathbf P_i\in\mathrm{range}\,\Pi_{\rm N}\cap\Hsp$, and Eq.~\eqref{eq:callaway} gives
$\Cop_{\rm Cal}\jj^{(i)}=\jj^{(i)}/\tau_R$. Then $K(p)=|\jj|^2/(p+\tau_R^{-1})$, and
$\lambda=|\jj|^2\tau_R/(k_BT^2)=\Lambda_i\tau_R=\rho c\,v^2\tau_R/3$ by Eq.~\eqref{eq:harmonic} and cubic
symmetry. The diagonal elements of the projectors are $e_\nu^2/|\ee|^2+\sum_iP_{i,\nu}^2/|\mathbf P_i|^2$
and $e_\nu^2/|\ee|^2$, both $O(n^{-1})$.
\end{proof}

Eliminating the deviatoric second moment of the same model under $|p|\tau_N\ll1$ and
$v\tau_N/L_{\rm v}\ll1$, with $L_{\rm v}$ the scale of spatial variation, gives the conserving GK
closure~\cite{GuoWang2022,Sendra2021,Sendra2022}
\begin{equation}
\tau_R\partial_t\mathbf q+\mathbf q=-\lambda\nabla T+\ell^2\big[\nabla^2\mathbf q+\beta\,\nabla(\nabla\cdot\mathbf q)\big],
\label{eq:gk}
\end{equation}
with $\beta=1/3$, $\ell^2=v^2\tau_c\tau_R/5$ and $\tau_c=(\tau_N^{-1}+\tau_R^{-1})^{-1}$; the original
GK equation has $\beta=2$ and $\tau_N$ in place of $\tau_c$~\cite{GuyerKrumhansl1966b}. Its relaxation
time is the $\tau_R$ of Proposition~\ref{prop:callaway}, i.e., $\tmem$. The mode lifetimes, however, are
close to $\tau_c$, not to $\tau_R$: lifetimes alone do not give the Cattaneo time even in the grey
model, whereas within this two-parameter family the pair $(\tau_c,\lambda)$ fixes $(\tau_N,\tau_R)$.
Section~\ref{sec:identifiability} asks what remains of this when the operator is not restricted to a
two-parameter family.

\subsection{Static data and operator classes}\label{sec:classes}

\begin{definition}[Static data]\label{def:data}
For $\Cop$ satisfying (K2)--(K3), the static data are $\mathcal D(\Cop)=(\mathbf r,\ker\Cop,K_0)$, where
$\mathbf r=\mathrm{diag}\,\Cop$ are the mode relaxation rates (lifetimes $\tau_\nu=1/r_\nu$). The
harmonic vectors $\ee$ and $\jj$ are fixed. For several directions, $K_0$ is replaced by the dc tensor.
\end{definition}

Under (K3) all $r_\nu>0$: a vanishing diagonal element of a positive-semidefinite matrix forces the
corresponding row to vanish, which would put a unit vector in $\ker\Cop=\mathrm{span}\{\ee\}$, whereas
$\ee$ has no zero component. The rates $r_\nu$ include the self-coupling of repeated slots in an event;
they are not, in general, the self-energy linewidths of RTA codes (see Lemma~\ref{lem:rta}).

\begin{definition}[Operator classes]\label{def:classes}
(P)~The positive-semidefinite class is
$\mathfrak P(\ee)=\{\Cop\in\mathbb R^{n\times n}:\Cop=\Cop\tr\succeq0,\ \ker\Cop=\mathrm{span}\{\ee\}\}$.
(E)~Let $Z=(\bm\zeta_\gamma)_{\gamma\in\mathcal E}$ be a fixed finite family of vectors in $\Hsp$ that
spans $\Hsp$. For rates $g\in[0,\infty)^{\mathcal E}$ put
\begin{equation}
\Cop(g)=\sum_{\gamma\in\mathcal E}g_\gamma\,\bm\zeta_\gamma\bm\zeta_\gamma\tr .
\label{eq:cone}
\end{equation}
The rates are admissible if $\ker\Cop(g)=\mathrm{span}\{\ee\}$, equivalently if
$\{\bm\zeta_\gamma:g_\gamma>0\}$ spans $\Hsp$. The fixed event cone is
$\mathfrak E(Z)=\{\Cop(g):g\ \text{admissible}\}\subset\mathfrak P(\ee)$.
\end{definition}

Three-phonon kinetics has the form~\eqref{eq:cone}~\cite{Fugallo2013}. For an event
$\gamma:\ \mathsf p\to\mathsf a+\mathsf b$ allowed by crystal-momentum conservation (normal or umklapp)
and by energy conservation on the mesh, let $\mathbf s_\gamma=-\bm\delta_{\mathsf p}+\bm\delta_{\mathsf a}+\bm\delta_{\mathsf b}$
be its stoichiometric vector ($-\bm\delta_{\mathsf p}+2\bm\delta_{\mathsf a}$ for a repeated daughter).
Linearizing the net rate $G_\gamma[f_{\mathsf p}(1+f_{\mathsf a})(1+f_{\mathsf b})-(1+f_{\mathsf p})f_{\mathsf a}f_{\mathsf b}]$
about $f^0$ and using detailed balance gives $-G_\gamma P_\gamma\,\mathbf s_\gamma\tr\mathsf W^{-2}\delta\mathbf f$,
with $P_\gamma=f^0_{\mathsf p}(1+f^0_{\mathsf a})(1+f^0_{\mathsf b})$, so that Eq.~\eqref{eq:kinetic} holds
with
\begin{equation}
\bm\zeta_\gamma=\mathsf W^{-1}\mathbf s_\gamma,\qquad g_\gamma=G_\gamma P_\gamma\ge0 .
\label{eq:events}
\end{equation}
Energy conservation $\mathbf s_\gamma\tr\bm\varepsilon=0$ gives $\bm\zeta_\gamma\in\Hsp$. The event set,
hence $Z$, is fixed by the harmonic spectrum; the rates $g_\gamma$ are fixed by the anharmonic
amplitudes, which are the unknowns of the inverse problem considered here. On a finite mesh, exact
resonance is exceptional; with a broadened resonance condition the accepted events must be given an
energy-conserving stoichiometry, which is a declared approximation (Supplemental Material, Sec.~S1).
The class statements below concern any fixed finite $Z$.
For given data $(\mathbf r^\star,K^\star)$ the feasible set is
\begin{equation}
\begin{split}
\mathcal F(\mathbf r^\star,K^\star)=\{g\ \text{admissible}:\ &\mathrm{diag}\,\Cop(g)=\mathbf r^\star,\\
&K_0(g)=K^\star\}.
\end{split}
\label{eq:feasible}
\end{equation}

Compressing the kinetics onto fewer variables is a separate question. Block observables
$\mathbf N=\mathsf B\yy$ (for instance the populations of degenerate blocks) obey a closed equation
$\dot{\mathbf N}=-\mathsf G\mathbf N$ for every initial state if and only if $\mathsf B\Cop=\mathsf G\mathsf B$,
the exact-lumpability condition~\cite{WeiKuo1969,Aoki1968,LiRabitz1989,Buchholz1994}; otherwise the reduced
dynamics contains memory terms of projection type~\cite{Zwanzig1961,Mori1965}. We use this known
criterion only as background.

% Section III — Static data do not determine the memory time
\section{Static data do not determine the memory time}\label{sec:identifiability}

Fix the harmonic vectors $\ee$ and $\jj\in\Hsp\setminus\{0\}$ and an operator class
$\mathfrak X\in\{\mathfrak P(\ee),\mathfrak E(Z)\}$. The question of this section is the size of the set
$\{\tmem(\Cop):\Cop\in\mathfrak X,\ \mathcal D(\Cop)=\mathcal D^\star\}$ for given static data
$\mathcal D^\star$. In Example~\ref{ex:four} and Theorem~\ref{thm:psd} rates are measured in units of a
common rate and times in its inverse.

\subsection{The positive-semidefinite class}\label{sec:psd}

\begin{example}[Same data, different memory; \status{proved}]\label{ex:four}
Let $n=4$, $\ee=(1,1,1,1)\tr/2$ and $\jj=(1,0,0,-1)\tr$. Let $\Cop_A=\tfrac43I-\tfrac13\mathbf 1\mathbf 1\tr$
and $\Cop_B=I-\tfrac12\mathsf A_4$, where $\mathbf 1=(1,1,1,1)\tr$ and $\mathsf A_4$ is the adjacency
matrix of the cycle $1\!-\!2\!-\!3\!-\!4\!-\!1$. Both operators have unit diagonal and kernel
$\mathrm{span}\{\ee\}$, and
\begin{equation}
K_A(p)=\frac{2}{p+4/3},\qquad K_B(p)=\frac1{p+1}+\frac1{p+2}.
\label{eq:four}
\end{equation}
Hence $K_A(0)=K_B(0)=3/2$, whereas $\tmem=3/4$ for $\Cop_A$ and $5/6$ for $\Cop_B$.
\end{example}

Both operators are weighted graph Laplacians on the four modes, so both lie in the single event cone
generated by the six vectors $\bm\delta_\mu-\bm\delta_\nu$. Static data therefore determine neither
the operator nor $\tmem$ even inside a fixed cone; Theorem~\ref{thm:cone} below shows that they bound
$\tmem$ there. In the positive-semidefinite class no bound exists.

\begin{theorem}[Unbounded memory at fixed static data; \status{proved}]\label{thm:psd}
Let $n=8$ and label the modes by $i\in\{0,\dots,7\}$. Let $\mathsf H$ be the symmetric orthogonal
Walsh matrix $\mathsf H_{ik}=8^{-1/2}(-1)^{\langle i,k\rangle}$, where $\langle i,k\rangle$ counts the
binary digits equal to one in both $i$ and $k$, and let $\mathbf h_k$ be its columns. For $0<\eta<1$ let
$\mathsf T(\eta)$ be the symmetric matrix whose only nonzero entries are
\begin{equation}
\begin{gathered}
\mathrm{diag}\,\mathsf T=(0,\eta^2,2,2-\eta^2,1,1,1,1),\\
\mathsf T_{12}=\mathsf T_{21}=\eta,\qquad \mathsf T_{47}=\mathsf T_{74}=-\eta ,
\end{gathered}
\label{eq:walsh}
\end{equation}
and put $\Cop_\eta=\mathsf H\,\mathsf T(\eta)\,\mathsf H\tr$, $\ee=\mathbf h_0$ and $\jj=\mathbf h_2$. Then:
\begin{enumerate}
\item[(i)] $\Cop_\eta\in\mathfrak P(\ee)$;
\item[(ii)] $\mathrm{diag}\,\Cop_\eta=(1,\dots,1)$;
\item[(iii)] $K_\eta(p)=(p+\eta^2)/[p^2+(2+\eta^2)p+\eta^2]$, hence $K_\eta(0)=1$ and
$\tmem(\Cop_\eta)=1+\eta^{-2}$;
\item[(iv)] the fixed-point-free involution $i\mapsto i\oplus3$ (bitwise exclusive or) commutes with
$\Cop_\eta$, leaves $\ee$ invariant and reverses $\jj$;
\item[(v)] for $|\omega|\ge\omega_{\min}>0$, $\big|K_\eta(i\omega)\,(i\omega+2)-1\big|\le\eta^2/(2\omega_{\min})$.
\end{enumerate}
Consequently, in $\mathfrak P(\ee)$ the static data
$(\mathbf r,\ker\Cop,K_0)=\big((1,\dots,1),\mathrm{span}\{\ee\},1\big)$ are compatible with every value
$\tmem\in(2,\infty)$.
\end{theorem}

\begin{proof}
(i) $\mathsf H$ is the character table of $(\mathbb Z_2)^3$ normalized to be orthogonal; it is symmetric.
$\mathsf T(\eta)$ is block diagonal on the index sets $\{0\}$, $\{1,2\}$, $\{3\}$, $\{4,7\}$, $\{5\}$,
$\{6\}$. The block on $\{1,2\}$ has determinant $\eta^2>0$ and positive trace; the block on $\{4,7\}$ has
determinant $1-\eta^2>0$; the remaining diagonal entries are $0$, $2-\eta^2>0$, $1$, $1$. Hence
$\mathsf T\succeq0$ with kernel $\mathrm{span}\{\bm\delta_0\}$, and $\Cop_\eta\succeq0$ with kernel
$\mathrm{span}\{\mathsf H\bm\delta_0\}=\mathrm{span}\{\ee\}$.
(ii) Since $\mathsf H_{ik}^2=1/8$ and $\mathrm{tr}\,\mathsf T=8$,
$(\Cop_\eta)_{ii}=1+2\eta(\mathsf H_{i1}\mathsf H_{i2}-\mathsf H_{i4}\mathsf H_{i7})$. Parity of
$\langle i,\cdot\rangle$ is additive under exclusive or, so
$\mathsf H_{i1}\mathsf H_{i2}=\tfrac18(-1)^{\langle i,3\rangle}=\mathsf H_{i4}\mathsf H_{i7}$, and the
bracket vanishes.
(iii) $(p+\Cop_\eta)^{-1}=\mathsf H(p+\mathsf T)^{-1}\mathsf H\tr$, so $K_\eta(p)$ is the $(2,2)$ entry of
the inverse of $p+\mathsf T$ restricted to $\{1,2\}$, which is $(p+\eta^2)/[(p+\eta^2)(p+2)-\eta^2]$. At
$p=0$ this is $1$; differentiating, $-K_\eta'(0)=1+\eta^{-2}$. Equivalently,
$\Cop_\eta(-\mathbf h_1/\eta+\mathbf h_2)=\jj$, so $\xx=-\mathbf h_1/\eta+\mathbf h_2$, $K_0=1$ and
$N=1+\eta^{-2}$.
(iv) Let $\mathsf P$ be the permutation $i\mapsto i\oplus3$. Then $\mathsf P\mathbf h_k=(-1)^{\langle3,k\rangle}\mathbf h_k$,
so $\mathsf P$ commutes with $\Cop_\eta$ if and only if $\mathsf T$ couples only indices of equal sign
$(-1)^{\langle 3,k\rangle}$. This holds: the signs are $-1$ for $k=1,2$ and $+1$ for $k=4,7$. The sign
is $+1$ for $k=0$ and $-1$ for $k=2$.
(v) $K_\eta(p)(p+2)-1=\eta^2/D(p)$ with $D(p)=p^2+(2+\eta^2)p+\eta^2$, and
$|D(i\omega)|\ge|\mathrm{Im}\,D(i\omega)|=(2+\eta^2)|\omega|\ge2|\omega|$.
For the consequence, take $\eta=(\tmem-1)^{-1/2}$.
\end{proof}

\begin{remark}[Mechanism]
$D$ has a root $-\eta^2/2+O(\eta^4)$ whose residue in $K_\eta$ is $\eta^2/4+O(\eta^4)$: a relaxation mode
of vanishing current weight and diverging lifetime carries about half of $K_0$ and all of the divergence
of $N$. Part~(v) shows that at frequencies $|\omega|\ge\omega_{\min}$ the response is within a relative
$\eta^2/(2\omega_{\min})$ of the single fast pole $1/(p+2)$. A band-limited measurement does not see
the divergence of $\tmem$; the slope at $p=0$ is not a property of the response in a finite band.
\end{remark}

\begin{remark}[Spectral reading and attribution]\label{rem:stieltjes}
By Eq.~\eqref{eq:stieltjes}, $|\jj|^2$, $K_0$ and $N$ are the moments of orders $0$, $-1$ and $-2$ of
the spectral measure of $\Cop|_{\Hsp}$ in the state $\jj$. The static data and the harmonic spectrum fix
the first two, and $\tmem$ requires the third. Without further structure no bound follows: this is the
elementary indeterminacy of a Stieltjes function from finitely many of its moments~\cite{Akhiezer2020,Schmudgen2017};
the same representation underlies classical bounds on effective transport coefficients of
composites~\cite{Golden1983}.
A spectral gap $\Gamma_{\min}$ gives $\tmem\le\Gamma_{\min}^{-1}$ (Lemma~\ref{lem:forms}). The content
of Theorem~\ref{thm:psd} is that prescribing the diagonal and the kernel does not supply such
structure; after a diagonal scaling, the positive-semidefinite matrices with prescribed diagonal form an
affine image of the elliptope~\cite{Saunderson2012}. The physical expectation is anticipated by the parity
structure of relaxons~\cite{Simoncelli2020} and by numerical failures of rate-only
descriptions~\cite{MalviyaRavichandran2023}; we have not found the identifiability statement itself for
phonon kinetics in the literature searched.
\end{remark}

\begin{remark}[The family is not microscopic]\label{rem:notcone}
No fixed finite event cone contains $\Cop_\eta$ for all small $\eta$. Let $\mathbf u=\mathbf h_1$; then
$\mathbf u\tr\Cop_\eta\jj=\eta$ and $\mathbf u\tr\Cop_\eta\mathbf u=\eta^2$. For
$\Cop=\sum_\gamma g_\gamma\bm\zeta_\gamma\bm\zeta_\gamma\tr$ with $g\ge0$,
\begin{equation}
|\mathbf u\tr\Cop\jj|\le\sum_\gamma g_\gamma|\mathbf u\tr\bm\zeta_\gamma|\,|\jj\tr\bm\zeta_\gamma|\le R\,\mathbf u\tr\Cop\mathbf u,
\label{eq:obstruction}
\end{equation}
with $R=\max\{|\jj\tr\bm\zeta_\gamma|/|\mathbf u\tr\bm\zeta_\gamma|:\mathbf u\tr\bm\zeta_\gamma\ne0\}$
(events with $\mathbf u\tr\bm\zeta_\gamma=0$ do not contribute). Membership for all small $\eta$ would
require $\eta\le R\eta^2$, which fails for $\eta<1/R$; if $R=0$ or no event has
$\mathbf u\tr\bm\zeta_\gamma\ne0$, the cross term vanishes, which contradicts $\eta>0$.
\end{remark}

\subsection{What static data do determine}\label{sec:lower}

\begin{proposition}[Lower bound from static data; \status{proved}]\label{prop:lower}
For every $\Cop\in\mathfrak P(\ee)$ and $\jj\in\Hsp\setminus\{0\}$,
\begin{equation}
\tmem\ \ge\ \frac{|\xx|}{|\jj|}\ \ge\ \tstat=\frac{K_0}{|\jj|^2},
\label{eq:lower}
\end{equation}
and $\tmem=\tstat$ if and only if $\jj$ is an eigenvector of $\Cop$.
\end{proposition}

\begin{proof}
By the Cauchy--Schwarz inequality $K_0=\jj\tr\xx\le|\jj|\,|\xx|$, so $|\xx|/K_0\ge1/|\jj|$ and
$\tmem=|\xx|\,(|\xx|/K_0)\ge|\xx|/|\jj|\ge K_0/|\jj|^2$. Equality throughout holds if and only if
$\xx\parallel\jj$, that is, $\Cop\jj\parallel\jj$.
\end{proof}

By Eq.~\eqref{eq:harmonic}, $\tstat=\lambda/\Lambda_i$ is the time of the elementary kinetic estimate
$\lambda=\Lambda_i\tau$; it is computable from the dc conductivity and the harmonic spectrum. In the
RTA it is the $c_\nu v_{\nu,i}^2$-weighted mean lifetime, and Proposition~\ref{prop:lower} reduces to
$\tmem^{\rm RTA}\ge\tstat^{\rm RTA}$. Equality holds in the grey Callaway model
(Proposition~\ref{prop:callaway}). Static data thus fix a lower end of the range of $\tmem$; whether that
end is approached in a given event geometry is a numerical question \PENDING{cross-reference: minima of
$\tmem/\tstat$ found in the Debye and AlN event geometries, Sec.~\ref{sec:aln}}.

The static data are not independent of each other. For an event operator, the dc response is bounded
below by the RTA response built from the same diagonal, up to the factor $1/\varsigma_{\max}$.

\begin{lemma}[Consistency of lifetimes and dc response; \status{proved}]\label{lem:rta}
Let $g$ be admissible, $\mathbf r=\mathrm{diag}\,\Cop(g)$ and $K_{\rm RTA}=\sum_\nu j_\nu^2/r_\nu$, and let
$\varsigma_{\max}$ be the largest number of distinct modes touched by an active event ($3$ for
$\mathsf p\to\mathsf a+\mathsf b$ with $\mathsf a\ne\mathsf b$, $2$ for $\mathsf p\to2\mathsf a$, up to $4$ for
four-phonon events). Then $K_0\ge K_{\rm RTA}/\varsigma_{\max}$. Equality holds for some
$\jj\in\Hsp\setminus\{0\}$ if and only if $\Cop(g)\yy=\varsigma_{\max}\,\mathrm{diag}(\mathbf r)\,\yy$ for some
$\yy\ne0$, and the currents attaining it are $\jj=\Cop(g)\yy$ with such $\yy$; this requires every active
event to touch exactly $\varsigma_{\max}$ modes. The constant is attained for $\varsigma_{\max}=2$, $3$ and $4$.
When some active event touches fewer than $\varsigma_{\max}$ modes, as for three-phonon events together with
four-phonon events on four modes, the inequality is strict, but the constant cannot be improved.
\end{lemma}

\begin{proof}
For $\jj\in\mathrm{range}\,\Cop$, $K_0=\max_{\yy}(2\jj\tr\yy-\yy\tr\Cop\yy)$, the variational principle of
the linearized Boltzmann equation~\cite{Ziman2001}. For a vector $\bm\zeta$ with $|\mathrm{supp}\,\bm\zeta|=s$,
the Cauchy--Schwarz inequality gives $(\bm\zeta\tr\yy)^2\le s\sum_\nu\zeta_\nu^2y_\nu^2$, that is
$\bm\zeta\bm\zeta\tr\preceq s\,\mathrm{diag}(\bm\zeta\circ\bm\zeta)$; summing with the rates,
$\Cop\preceq\varsigma_{\max}\,\mathrm{diag}\,\Cop$. With $\yy=\mathrm{diag}(\mathbf r)^{-1}\jj/\varsigma_{\max}$,
$K_0\ge2K_{\rm RTA}/\varsigma_{\max}-K_{\rm RTA}/\varsigma_{\max}$. Equality forces
$\yy\tr(\varsigma_{\max}\,\mathrm{diag}\,\Cop-\Cop)\yy=0$ for this $\yy$, hence
$\Cop\yy=\varsigma_{\max}\,\mathrm{diag}(\mathbf r)\,\yy=\jj$, because the matrix is positive semidefinite.
Conversely, if $\Cop\yy=\varsigma_{\max}\,\mathrm{diag}(\mathbf r)\,\yy$ with $\yy\ne0$, then
$\jj=\Cop\yy\in\Hsp$ is nonzero because all $r_\nu>0$, and
$K_0=\yy\tr\Cop\yy=\varsigma_{\max}\,\yy\tr\mathrm{diag}(\mathbf r)\,\yy=K_{\rm RTA}/\varsigma_{\max}$. Each active event contributes to
$\yy\tr(\varsigma_{\max}\,\mathrm{diag}\,\Cop-\Cop)\yy$ a nonnegative term, which vanishes if and only if the
numbers $\zeta_{\gamma\nu}y_\nu$, $\nu\in\mathrm{supp}\,\bm\zeta_\gamma$, are equal, and equal to zero when
$|\mathrm{supp}\,\bm\zeta_\gamma|<\varsigma_{\max}$. A zero of $\yy$ on one mode of an active event therefore
spreads to all its modes. The active events of an admissible network connect all modes, since a set of
modes not connected to the others would carry an energy invariant of its own; so one active event with
fewer than $\varsigma_{\max}$ modes forces $\yy=0$. Exact networks attain the ratio $1/\varsigma_{\max}$ for
$\varsigma_{\max}=2$, $3$ and $4$, and a family of mixed networks approaches it (Supplemental Material,
Sec.~S2).
\end{proof}

The lemma refers to $\mathrm{diag}\,\Cop$, in which the self-coupling of a repeated daughter is counted
in full. A self-energy linewidth that counts it with half weight is a different quantity, and the
inequality can fail for it (an exact counterexample is given in the Supplemental Material, Sec.~S2).

\subsection{Fixed event cone}\label{sec:cone}

\begin{definition}[Basic solutions and the constant $M$]\label{def:M}
For $I\subset\mathcal E$ with $|I|=n_{\Hsp}-1$, consider the linear system
\begin{equation}
\bm\zeta_\gamma\tr\yy=0\ \ (\gamma\in I),\qquad \yy\in\Hsp,\qquad \jj\tr\yy=1 .
\label{eq:basic}
\end{equation}
It has a unique solution $\yy_I$ if and only if $\{\bm\zeta_\gamma\}_{\gamma\in I}\cup\{\jj\}$ is a basis
of $\Hsp$; such $I$ are called bases, and $\mathfrak B$ denotes their set, which is nonempty because $Z$
spans $\Hsp$ and $\jj\ne0$. Put $M=M(Z,\jj)=\max_{I\in\mathfrak B}|\yy_I|$.
\end{definition}

$M$ does not depend on the rates. It depends on $Z$ only through the hyperplanes
$\mathrm{span}\{\bm\zeta_\gamma\}_{\gamma\in I}$, and $M(Z,s\jj)=M(Z,\jj)/|s|$, so that $M^2K_0$ and
$M^2|\jj|^2$ do not depend on the normalization of $\jj$.

\begin{theorem}[Cone bound; \status{proved}]\label{thm:cone}
Let $Z$ span $\Hsp$, $\jj\in\Hsp\setminus\{0\}$ and $g$ admissible. Fix an orthonormal basis of $\Hsp$;
for every $I\subset\mathcal E$ with $|I|=n_{\Hsp}-1$ define $\bm\psi_I\in\Hsp$ by $\bm\psi_I\tr\yy=\det[\bm\zeta_I,\yy]$ for all
$\yy\in\Hsp$, the determinant having the columns $\bm\zeta_\gamma$ ($\gamma\in I$, in a fixed order)
followed by $\yy$, and put $g_I=\prod_{\gamma\in I}g_\gamma$. Then
\begin{equation}
\frac{\xx}{K_0}=\sum_{I\in\mathfrak B}\varpi_I\,\yy_I,\qquad
\varpi_I=\frac{g_I\,(\bm\psi_I\tr\jj)^2}{\sum_{J\in\mathfrak B}g_J\,(\bm\psi_J\tr\jj)^2},
\label{eq:hull}
\end{equation}
a convex combination over the active bases ($g_I>0$). Consequently
\begin{equation}
|\Cop(g)^+\jj|\le M\,K_0(g),\qquad \tmem(g)\le M^2K_0(g).
\label{eq:conebound}
\end{equation}
\end{theorem}

\begin{proof}
Identify $\Hsp$ with $\mathbb R^{n_{\Hsp}}$ through the orthonormal basis and let
$A=\Cop(g)|_{\Hsp}=Z\,\mathrm{diag}(g)\,Z\tr$, where $Z$ also denotes the matrix with columns
$\bm\zeta_\gamma$; $A$ is positive definite. For $\yy,\mathbf z\in\Hsp$ the matrix determinant lemma gives
$\det(A+\yy\mathbf z\tr)=\det A+\mathbf z\tr\mathrm{adj}(A)\,\yy$. Writing
$A+\yy\mathbf z\tr=[Z,\yy]\,\mathrm{diag}(g,1)\,[Z,\mathbf z]\tr$ and expanding by the Cauchy--Binet formula
over the $n_{\Hsp}$-element subsets of columns, the subsets without the last column give $\det A$ and
those containing it give $\sum_{|I|=n_{\Hsp}-1}g_I\det[\bm\zeta_I,\yy]\det[\bm\zeta_I,\mathbf z]$. Hence
\begin{equation}
\mathrm{adj}(A)=\sum_{|I|=n_{\Hsp}-1}g_I\,\bm\psi_I\bm\psi_I\tr ,
\label{eq:adj}
\end{equation}
where dependent subsets give $\bm\psi_I=0$. With $A^{-1}=\mathrm{adj}(A)/\det A$,
$\xx=\sum_Ig_I\bm\psi_I(\bm\psi_I\tr\jj)/\det A$ and $K_0=\sum_Ig_I(\bm\psi_I\tr\jj)^2/\det A>0$. Now
$\bm\psi_I\tr\jj=\det[\bm\zeta_I,\jj]\ne0$ exactly when $I\in\mathfrak B$, and then
$\bm\psi_I/(\bm\psi_I\tr\jj)$ satisfies Eq.~\eqref{eq:basic}, because $\bm\psi_I\tr\bm\zeta_\gamma$ is a
determinant with a repeated column for $\gamma\in I$; it is therefore $\yy_I$. Dividing $\xx$ by $K_0$
gives Eq.~\eqref{eq:hull}. Finally $|\xx|/K_0\le\sum_I\varpi_I|\yy_I|\le M$ and
$\tmem=(|\xx|/K_0)^2K_0$.
\end{proof}

\begin{remark}[Relation to weighted least squares]\label{rem:bt}
By Lagrange's rule, $\xx/K_0$ minimizes $\sum_\gamma g_\gamma(\bm\zeta_\gamma\tr\yy)^2$ over $\yy\in\Hsp$
subject to $\jj\tr\yy=1$: it is the solution of a weighted least-squares problem with one exact
constraint. Ben-Tal and Teboulle proved, by the same Cauchy--Binet argument, that the solution of a
weighted least-squares problem lies in the convex hull of the solutions of its square subsystems for
every positive weight vector~\cite[Theorem~2.1 and Corollary~2.1]{BenTalTeboulle1990} (see also
Ref.~\cite{Berg1986}). Theorem~\ref{thm:cone} is that statement for the constrained problem, which is the
limit of infinite weight on the row $\jj\tr\yy=1$. Weight-independent bounds on scaled projections and
pseudoinverses are classical~\cite{Stewart1989,OLeary1990,Todd1990,HankeNeumann1993} and govern the
complexity of interior-point methods~\cite{VavasisYe1996}; approximating a closely related
weight-uniform condition measure is NP-hard in general~\cite{Tuncel1999}, so for large event sets $M$
must be bounded by exploiting structure. We claim no novelty for the inequality; its use to bound the
memory moment of a kinetic operator by its dc response is potentially novel within the literature
searched.
\end{remark}

The bound extends to rates that are not admissible but satisfy $\jj\in\mathrm{range}\,\Cop(g)$. Put
$g_\theta=g+\theta\mathbf 1$ ($\theta>0$), which is admissible, and $E=\Cop(\mathbf 1)$. Split
$\Hsp=\mathcal R\oplus\mathcal Q$ with $\mathcal R=\mathrm{range}\,\Cop(g)$ and
$\mathcal Q=\ker\Cop(g)\cap\Hsp$, and write $\xx_\theta=\Cop(g_\theta)^+\jj=\mathbf w_\theta+\mathbf u_\theta$
accordingly. The $\mathcal Q$ component of $\Cop(g_\theta)\xx_\theta=\jj$ gives
$\mathbf u_\theta=-(\Pi_{\mathcal Q}E\Pi_{\mathcal Q})^{-1}\Pi_{\mathcal Q}E\,\mathbf w_\theta$, and the
$\mathcal R$ component shows $\mathbf w_\theta\to\mathbf w=\Cop(g)^+\jj$ as $\theta\to0$; hence
$\mathbf u_\theta\to\mathbf u_\infty$ and $K_0(g_\theta)=\jj\tr\mathbf w_\theta\to K_0(g)$. Since
$\mathbf w\perp\mathbf u_\infty$, $|\mathbf w|\le|\mathbf w+\mathbf u_\infty|\le M K_0(g)$ by
Theorem~\ref{thm:cone} and continuity.

\begin{proposition}[Sharpness without lifetime constraints; \status{proved}]\label{prop:sharp}
For every $K^\star>0$,
\begin{equation}
\sup\{\tmem(g):\ g\ \text{admissible},\ K_0(g)=K^\star\}=M^2K^\star .
\label{eq:sharp}
\end{equation}
The supremum is approached by rates concentrated on a basis $I^\ast$ with $|\yy_{I^\ast}|=M$ and is in
general not attained; along the approaching sequence the rates on $I^\ast$ diverge, and so do the
relaxation rates of the modes touched by $I^\ast$.
\end{proposition}

\begin{proof}
For $0<\epsilon<1$ let $g^\epsilon_\gamma=1$ on $I^\ast$ and $\epsilon$ elsewhere; $g^\epsilon$ is
admissible. For $I\in\mathfrak B$, $g^\epsilon_I=\epsilon^{|I\setminus I^\ast|}$ with
$|I\setminus I^\ast|\ge1$ for $I\ne I^\ast$, so $\varpi_{I^\ast}\to1$ in Eq.~\eqref{eq:hull} and
$|\xx|/K_0\to M$ as $\epsilon\to0$. Since $\Cop(sg)=s\Cop(g)$, the ratio $|\xx|/K_0$ is invariant under
$g\mapsto sg$ while $K_0$ scales as $1/s$; rescaling $g^\epsilon$ to $K_0=K^\star$ gives
$\tmem=(|\xx|/K_0)^2K^\star\to M^2K^\star$, and Theorem~\ref{thm:cone} gives the reverse inequality.
Every $n_{\Hsp}$-element subset of events contains an event outside $I^\ast$, and some $I^\ast\cup\{\gamma\}$
with $\gamma\notin I^\ast$ is a basis of $\Hsp$, because the $\bm\zeta_\gamma$ ($\gamma\in I^\ast$) are
independent and $Z$ spans $\Hsp$; so $\det A=\Theta(\epsilon)$ while
$\sum_Ig^\epsilon_I(\bm\psi_I\tr\jj)^2\to(\bm\psi_{I^\ast}\tr\jj)^2>0$; thus $K_0(g^\epsilon)\sim c/\epsilon$
and the rescaling factor $K_0(g^\epsilon)/K^\star$ diverges.
\end{proof}

\begin{remark}[With lifetimes fixed]\label{rem:lifetimes}
Since $\mathcal F(\mathbf r^\star,K^\star)$ is contained in the set of Eq.~\eqref{eq:sharp},
$\sup_{\mathcal F}\tmem\le M^2K^\star$. Equality is possible: if $\jj=\bm\zeta_\beta$ is parallel to an
event vector, no subset containing $\beta$ is a basis, $M$ is attained by some $J^\ast\not\ni\beta$, and
every $g$ supported on the basis $J^\ast\cup\{\beta\}$ of $\Hsp$ gives $\xx/K_0=\yy_{J^\ast}$ exactly; with
$\mathbf r^\star=\mathrm{diag}\,\Cop(g)$ and $K^\star=K_0(g)$ the lifetime-constrained supremum is then
attained (an exact example with seven modes and ten events is given in the Supplemental Material,
Sec.~S3). Physical currents~\eqref{eq:current} vanish where $v_{\nu,i}=0$; they are odd under time
reversal, so their support consists of pairs $\{\nu,\bar\nu\}$ of modes of equal energy, and they are never
parallel to a three-phonon event vector, whose support has three modes or two modes of different energies.
For them it is \status{open} whether $\sup_{\mathcal F}\tmem<M^2K^\star$. For rates tied by a symmetry
group, the proof of Proposition~\ref{prop:sharp} transfers when some maximizing $I^\ast$ is a union of
orbits; otherwise the supremum over symmetric rates is not known.
\end{remark}

For three-phonon events, $M$ has a reading in terms of sub-networks.

\begin{proposition}[Sub-network form of $M$; \status{proved}]\label{prop:subnetwork}
Let $\bm\zeta_\gamma=\mathsf W^{-1}\mathbf s_\gamma$ as in Eq.~\eqref{eq:events}, and for $I\in\mathfrak B$
let $\mathsf S_I$ be the matrix with columns $\mathbf s_\gamma$, $\gamma\in I$. Then
$\ker\mathsf S_I\tr=\mathrm{span}\{\bm\varepsilon,\bm\phi_I\}$ for some $\bm\phi_I\notin\mathrm{span}\{\bm\varepsilon\}$, and
\begin{equation}
\yy_I=\frac{\Pi_{\Hsp}\mathsf W\bm\phi_I}{\jj\tr\mathsf W\bm\phi_I},\qquad
|\yy_I|=\frac{|\Pi_{\Hsp}\mathsf W\bm\phi_I|}{|\jj\tr\mathsf W\bm\phi_I|}.
\label{eq:subnetwork}
\end{equation}
\end{proposition}

\begin{proof}
Since $\mathsf W$ is invertible, the $\bm\zeta_\gamma$ ($\gamma\in I$) are independent if and only if
$\mathsf S_I$ has full column rank $n-2$; then $\dim\ker\mathsf S_I\tr=2$, and $\bm\varepsilon\in\ker\mathsf S_I\tr$
by energy conservation. The conditions $\bm\zeta_\gamma\tr\yy=\mathbf s_\gamma\tr\mathsf W^{-1}\yy=0$
($\gamma\in I$) say $\mathsf W^{-1}\yy\in\ker\mathsf S_I\tr$, i.e.,
$\yy=a\,\mathcal V^{1/2}\ee+b\,\mathsf W\bm\phi_I$. The condition $\yy\in\Hsp$ gives
$\yy=b\,\Pi_{\Hsp}\mathsf W\bm\phi_I$, and $\jj\tr\yy=1$ with $\jj\perp\ee$ gives
$b=1/(\jj\tr\mathsf W\bm\phi_I)$; the denominator is nonzero because $I\in\mathfrak B$.
\end{proof}

The vector $\bm\phi_I$ is a conserved quantity of the events in $I$ other than energy,
$\mathbf s_\gamma\tr\bm\phi_I=0$ for $\gamma\in I$. $M$ is large when some $(n-2)$-event sub-network
has such a hidden invariant with a small overlap with the current. By Eq.~\eqref{eq:hull},
$|\xx|/K_0\le\sum_I\varpi_I|\yy_I|$, so a large ratio $\tmem/K_0$ requires convex weight on bases with
large $|\yy_I|$: for $0<\epsilon\le1/2$, $\tmem\ge(1-\epsilon)^2M^2K_0$ puts at least half of the weight on
bases with $|\yy_I|\ge(1-2\epsilon)M$. The weight resides in the products $g_I(\bm\psi_I\tr\jj)^2$, not
necessarily in individual rates. Proposition~\ref{prop:sharp} shows that concentrating the rates on one
such basis attains the bound, which turns the invariant of the sub-network into a slow, weakly
current-coupled relaxation mode, the mechanism of Theorem~\ref{thm:psd}.

\begin{corollary}[Width of the feasible range; \status{proved}]\label{cor:gap}
If $\mathcal F(\mathbf r^\star,K^\star)$ is nonempty, then
$\sup_{\mathcal F}\tmem\,/\,\inf_{\mathcal F}\tmem\le M^2|\jj|^2$.
\end{corollary}

\begin{proof}
Combine $\tmem\le M^2K^\star$ (Theorem~\ref{thm:cone}) with $\tmem\ge K^\star/|\jj|^2$
(Proposition~\ref{prop:lower}).
\end{proof}

The structure of the feasible set is convex in part. With the extended value
$K_0(g)=\sup_{\yy}\{2\jj\tr\yy-\yy\tr\Cop(g)\yy\}$, equal to $\jj\tr\Cop(g)^+\jj$ if
$\jj\in\mathrm{range}\,\Cop(g)$ and $+\infty$ otherwise, $K_0$ is convex, lower semicontinuous and
nonincreasing in $g$, with $\partial K_0/\partial g_\gamma=-(\bm\zeta_\gamma\tr\xx)^2$ at admissible $g$.
By the generalized Schur complement,
\begin{multline}
\{g\ge0:\ \mathrm{diag}\,\Cop(g)=\mathbf r^\star,\ K_0(g)\le K^\star\}\\
=\Big\{g\ge0:\ \mathrm{diag}\,\Cop(g)=\mathbf r^\star,\
\begin{pmatrix}\Cop(g)&\jj\\ \jj\tr&K^\star\end{pmatrix}\succeq0\Big\},
\label{eq:lmi}
\end{multline}
so the infimum of $K_0$ at fixed lifetimes is a semidefinite program (\status{proved}). The set contains
nonadmissible rates; the infimum over admissible rates coincides with it whenever the lifetime
constraint admits an admissible point. No convexity of $\tmem$ is claimed.

\subsection{What bounds the memory: amplitude information}\label{sec:prior}

\begin{proposition}[Per-event prior; \status{proved}]\label{prop:prior}
Let $g^{\rm ref}$ be admissible and $F>0$. Every $g$ with $g_\gamma\ge g^{\rm ref}_\gamma/F$ for all
$\gamma\in\mathcal E$ is admissible and satisfies
\begin{equation}
\tmem(g)\le\frac{F}{\Gamma_{\min}\big(\Cop(g^{\rm ref})\big)},
\label{eq:prior}
\end{equation}
where $\Gamma_{\min}$ is the smallest eigenvalue on $\Hsp$. No lifetime or dc data enter, and the bound
holds with or without symmetry constraints on the rates. Conversely, for $\tau>0$, $\tmem(g)>\tau$ requires
$\min\{g_\gamma/g^{\rm ref}_\gamma:g^{\rm ref}_\gamma>0\}<[\tau\,\Gamma_{\min}(\Cop(g^{\rm ref}))]^{-1}$.
\end{proposition}

\begin{proof}
$\Cop(g)-F^{-1}\Cop(g^{\rm ref})=\sum_\gamma(g_\gamma-g^{\rm ref}_\gamma/F)\bm\zeta_\gamma\bm\zeta_\gamma\tr\succeq0$,
so $\Gamma_{\min}(\Cop(g))\ge\Gamma_{\min}(\Cop(g^{\rm ref}))/F>0$, which implies admissibility, and
Lemma~\ref{lem:forms} gives $\tmem(g)\le1/\Gamma_{\min}(\Cop(g))$. For the converse, if $g_\gamma=0$ for
some $\gamma$ with $g^{\rm ref}_\gamma>0$, the minimum is $0$ and there is nothing to prove; otherwise
$g\ge g^{\rm ref}/F$ with $F=\max\{g^{\rm ref}_\gamma/g_\gamma:g^{\rm ref}_\gamma>0\}$, and the bound gives
the converse.
\end{proof}

The information that confines $\tmem$ is therefore microscopic: a statement about every event rate
relative to a reference, which is what an anharmonic amplitude model provides and what static data do
not. The quantitative consequences for AlN are reported in Sec.~\ref{sec:aln}.

\subsection{Infrared criterion}\label{sec:ir}

We now let the mesh be refined. Consider uniform $\Gamma$-centered meshes $\mathbf k=\mathsf B\mathbf n/N_1$
($\mathbf n\in\mathbb Z^d$, $\mathsf B$ the matrix of primitive reciprocal-lattice vectors, so that
$|\det\mathsf B|\,V_{\rm cell}=(2\pi)^d$) with $N_1$ points per reciprocal-lattice direction in spatial
dimension $d$ ($N_{\mathbf k}=N_1^d$), the $\Gamma$ point excluded. With the normalization~\eqref{eq:current}, mode sums $\sum_\nu F(\mathbf k_\nu,s_\nu)/\mathcal V$
are Riemann sums for $(2\pi)^{-d}\sum_s\int_{\rm BZ}F\,d^d\mathbf k$. Write $j_\nu=\mathcal V^{-1/2}J(\mathbf k_\nu,s_\nu)$
with $J=w\varepsilon v_i$, and define the RTA moments
\begin{equation}
\mathcal M_m=\sum_s\int_{\rm BZ}J(\mathbf k,s)^2\,r(\mathbf k,s)^{-m}\,\frac{d^d\mathbf k}{(2\pi)^d},\qquad m=0,1,2.
\label{eq:Mm}
\end{equation}
On the meshes, $|\jj|^2$, $K_{\rm RTA}$ and $\tmem^{\rm RTA}$ are $\mathcal M_0(N_1)$, $\mathcal M_1(N_1)$
and $\mathcal M_2(N_1)/\mathcal M_1(N_1)$, with the mesh moments
$\mathcal M_m(N_1)=\mathcal V^{-1}\sum_\nu J(\mathbf k_\nu,s_\nu)^2r_\nu^{-m}$; their continuum counterparts are
$\mathcal M_0$, $\mathcal M_1$ and $\mathcal M_2/\mathcal M_1$.
On each acoustic branch, with $k=|\mathbf k|$ and $\hat{\mathbf k}=\mathbf k/k$, the hypotheses are:
\begin{enumerate}
\item[(H1)] $r(\mathbf k)=r_0(\hat{\mathbf k})\,k^\alpha[1+o(1)]$ as $k\to0$, uniformly in $\hat{\mathbf k}$, with
$0<r_-\le r_0(\hat{\mathbf k})\le r_+<\infty$;
\item[(H2)] $J(\mathbf k)\to J_0(\hat{\mathbf k})$ as $k\to0$, uniformly in $\hat{\mathbf k}$; $J$ is bounded on the
Brillouin zone; and $0<A_m=\sum_{\rm ac}\int_{S^{d-1}}J_0^2\,r_0^{-m}\,d\Omega<\infty$ for the orders $m$ considered;
\item[(H3)] away from $\mathbf k=0$, and on optical branches, $r$ is bounded below by a positive constant.
\end{enumerate}
For a branch with linear dispersion and sound velocity $v_s$, $w\varepsilon\to k_BT$ and
$J_0^2=(k_BTv_s)^2\hat k_i^2$, which vanishes on the plane $\hat k_i=0$; (H2) requires only that the
angular integral be positive. For the event operator one more hypothesis is needed:
\begin{enumerate}
\item[(H4)] for the admissible operators $\Cop_{N_1}$ on the meshes and $\xx=\Cop_{N_1}^+\jj$, there are
$c>0$ and $k_0>0$, independent of $N_1$, such that $|x_\nu|\ge c\,|j_\nu|/r_\nu$ for every mode with
$k_\nu<k_0$.
\end{enumerate}
Limits on the meshes require, in addition, some regularity:
\begin{enumerate}
\item[(H5)] for the orders $m$ considered, $\Phi_m=\sum_{\rm ac}J_0^2r_0^{-m}$ is bounded on $S^{d-1}$ and
continuous outside a closed set of directions of zero measure, and on every branch $J^2r^{-m}$ is Riemann
integrable on the part $k\ge\delta$ of the Brillouin zone for every $\delta>0$.
\end{enumerate}
Hypothesis (H5) admits the discontinuities of frequency-sorted acoustic branches at the acoustic axes,
where they touch, and where they cross; some regularity of this kind is necessary (Supplemental Material,
Sec.~S4).

\begin{theorem}[Infrared criterion]\label{thm:ir}
Let the meshes, the moments $\mathcal M_m$ and hypotheses (H1)--(H5) be as above.
\begin{enumerate}
\item[(i)] \emph{RTA moments} (\status{proved}). Assume (H1)--(H3). Then $\mathcal M_m<\infty$ if and
only if $m\alpha<d$; for $m\alpha=d$ the moment restricted to $k>k_{\min}$ diverges like $\log(1/k_{\min})$.
In the RTA the conductivity is therefore finite if and only if $\alpha<d$, and the memory moment if and
only if $2\alpha<d$; for $\alpha\ge d$, where $\mathcal M_1$ and $\mathcal M_2$ both diverge, the ratio of
the moments restricted to $k>k_{\min}$ diverges like $k_{\min}^{-\alpha}$ for $\alpha>d$ and like
$k_{\min}^{-d}/\log(1/k_{\min})$ for $\alpha=d$. If, in addition, (H5) holds, then
$\mathcal M_m(N_1)\to\mathcal M_m$ for $m\alpha<d$,
$\mathcal M_m(N_1)=(2\pi)^{-d}A_m\log N_1+o(\log N_1)$ for $m\alpha=d$, and
$\mathcal M_m(N_1)=N_1^{m\alpha-d}S_m/V_{\rm cell}\,[1+o(1)]$ for $m\alpha>d$, where
$S_m=\sum_{\mathbf n\ne0}\Phi_m(\mathsf B\mathbf n/|\mathsf B\mathbf n|)\,|\mathsf B\mathbf n|^{-m\alpha}\in(0,\infty)$.
Hence on the meshes $\tmem^{\rm RTA}$ converges to $\mathcal M_2/\mathcal M_1$ for $2\alpha<d$ and grows
like $\log N_1$ for $2\alpha=d$, like $N_1^{2\alpha-d}$ for $d/2<\alpha<d$, and like $N_1^d/\log N_1$ for
$\alpha=d$.
\item[(ii)] \emph{Event operator} (\status{derived under assumptions}, namely (H4) and the uniform form of
(H1)). Suppose that (H2) holds, that $\Phi_2$ is continuous outside a closed set of directions of zero
measure, that $\mathbf r=\mathrm{diag}\,\Cop_{N_1}$ satisfies (H1) uniformly in $N_1$, that is, with $r_0$
independent of $N_1$ and the $o(1)$ uniform in $\hat{\mathbf k}$ and in $N_1$, and that (H4) holds. If
$2\alpha\ge d$, then $N=|\Cop_{N_1}^+\jj|^2\to\infty$ and $\tmem\ge N^{1/2}/|\jj|\to\infty$ as
$N_1\to\infty$, with no condition on $K_0$; for $2\alpha>d$, $\tmem\gtrsim N_1^{\alpha-d/2}$, and for
$2\alpha=d$, $\tmem\gtrsim(\log N_1)^{1/2}$. These are lower bounds; two-sided mesh exponents for the
event operator do not follow.
\item[(iii)] \emph{Cone constant} (\status{derived under assumptions}, namely (H4) and the uniform form
of (H1)). Under the hypotheses of (ii) and $2\alpha\ge d$, $M_{N_1}^2K_{0,N_1}\to\infty$, where $M_{N_1}$ is the constant of
Theorem~\ref{thm:cone} for the event family on the mesh; if, in addition, $K_{0,N_1}$ stays bounded,
$M_{N_1}\to\infty$.
\end{enumerate}
\end{theorem}

\begin{proof}
(i) By (H3) and the boundedness of $J$, the integrand of Eq.~\eqref{eq:Mm} is bounded outside a ball
$k<k_0$. Inside it, (H1)--(H2) give
$J^2r^{-m}=[J_0(\hat{\mathbf k})^2r_0(\hat{\mathbf k})^{-m}+\epsilon(\mathbf k)]\,k^{-m\alpha}$ with
$\sup_{\hat{\mathbf k}}|\epsilon|\to0$ as $k\to0$; the error is additive because $J_0$ may vanish on a
plane. In polar coordinates $d^d\mathbf k=k^{d-1}dk\,d\Omega$ the ball contributes
$\int_{k_{\min}}^{k_0}k^{d-1-m\alpha}[A+\bar\epsilon(k)]\,dk$, where $A=A_m>0$ and $\bar\epsilon(k)\to0$.
Choosing $k_0$ with $|\bar\epsilon|\le A/2$ bounds this between $A/2$ and $3A/2$ times
$\int_{k_{\min}}^{k_0}k^{d-1-m\alpha}dk$, which stays bounded as $k_{\min}\to0$ if and only if $m\alpha<d$,
equals $\log(k_0/k_{\min})$ for $m\alpha=d$ and is of order $k_{\min}^{d-m\alpha}$ for $m\alpha>d$; the
statement for $\alpha\ge d$ follows. For the mesh laws, the optical branches and the part $k\ge k_0$ of
the acoustic branches contribute a bounded amount to $\mathcal M_m(N_1)$, by (H3) and the boundedness of
$J$. On the acoustic
branches the part $k_\nu<k_0$ equals
$N_1^{m\alpha-d}V_{\rm cell}^{-1}\sum_{0<|\mathsf B\mathbf n|<k_0N_1}[\varphi(\mathbf n)+\epsilon_{\rm ac}(\mathsf B\mathbf n/N_1)\,|\mathsf B\mathbf n|^{-m\alpha}]$,
where $\varphi(\mathbf n)=\Phi_m(\mathsf B\mathbf n/|\mathsf B\mathbf n|)\,|\mathsf B\mathbf n|^{-m\alpha}$ and
$\epsilon_{\rm ac}$ is the sum of $\epsilon$ over the acoustic branches. Lemma~S1 of the Supplemental
Material (Sec.~S4) evaluates such sums. On each dyadic shell $2^l\le|\mathsf B\mathbf n|<2^{l+1}$ the sum
of $\varphi$ is $2^{l(d-m\alpha)}$ times a Riemann sum of the homogeneous extension of $\varphi$ over the
shell $1\le|\mathsf B\mathbf u|<2$, which converges to a positive limit by (H5); the $\epsilon_{\rm ac}$
terms are bounded shell by shell. Summing over the shells, with $|\det\mathsf B|\,V_{\rm cell}=(2\pi)^d$, gives
the stated laws for $m\alpha\ge d$. For $m\alpha<d$ the part $k<\delta$ of $\mathcal M_m(N_1)$ and of
$\mathcal M_m$ is $O(\delta^{d-m\alpha})$ uniformly in $N_1$, and the remaining sums are Riemann sums that
converge by (H5), so that $\mathcal M_m(N_1)\to\mathcal M_m$. Taking the ratio $\mathcal M_2(N_1)/\mathcal M_1(N_1)$ gives the four
cases.
(ii) By (H4), $N\ge\sum_{k_\nu<k_0}x_\nu^2\ge c^2\sum_{k_\nu<k_0}j_\nu^2r_\nu^{-2}$. Only the upper side
of the uniform (H1) is used: $r_\nu\le r_0(\hat{\mathbf k}_\nu)k_\nu^\alpha[1+o(1)]$ with an $o(1)$
independent of $N_1$, so the last sum is bounded below by a mesh sum of the form treated in part~(i),
with angular function $\Phi_2$ and an error independent of $N_1$; by Lemma~S1 it grows like
$N_1^{2\alpha-d}$ for $2\alpha>d$ and like $\log N_1$ for $2\alpha=d$. Since $|\jj|^2$ is at most the
number of branches times $\sup J^2/V_{\rm cell}$, Proposition~\ref{prop:lower} gives
$\tmem\ge N^{1/2}/|\jj|\to\infty$ at the stated rates.
(iii) Theorem~\ref{thm:cone} gives $\tmem\le M_{N_1}^2K_{0,N_1}$ on every mesh.
\end{proof}

\begin{example}[Finite dc response, infinite slope; \status{proved}]\label{ex:ir}
The Stieltjes function
$K(p)=\int_0^1q^2(q^2+p)^{-1}dq=1-\sqrt p\,\arctan(1/\sqrt p)$ is the RTA response of a one-dimensional
branch with $J^2\propto q^2$ and $r\propto q^2$. It has $K(0)=1$ and $K(p)=1-\tfrac\pi2\sqrt p+O(p)$: the
small-$p$ expansion is not analytic and no Cattaneo time exists. With a cutoff $q>1/N_1$,
$K_0=1-1/N_1$ and $\tmem=N_1$.
\end{example}

Hypothesis (H4) is a property of a given family of operators, not of a class: it states that the
long-wavelength part of the solution $\xx$ is RTA-like, and it holds with $c=1$ in the RTA. In three
dimensions a finite conductivity requires $\alpha<3$, whereas a finite memory moment requires
$\alpha<3/2$. The low-frequency three-phonon rates of acoustic phonons follow power laws whose exponents
depend on the branch and on the crystal symmetry~\cite{Herring1954}; whether the operators of a given
material satisfy (H1)--(H5), and with which $\alpha$, is a separate question \PENDING{AlN: infrared law
of the acoustic rates and status of (H1)--(H5) for the AlN operators, Sec.~\ref{sec:aln}}. The mesh laws
are approached slowly: in the lattice test of Supplemental Sec.~S4, $\mathcal M_1(N_1)$ converges like
$N_1^{-(d-\alpha)}$ for $d/2<\alpha<d$, and at $\alpha=d$ the relative correction decays like
$1/\log N_1$, so that exponents fitted on coarse meshes are biased. When the
memory moment diverges, a Cattaneo time exists only for a regularized problem: a finite sample, whose
boundaries scatter long-wavelength phonons, or a finite frequency band. This motivates the
film-regularized memory time of Sec.~\ref{sec:aln}.

% Section IV — Model of the front-face measurement
\section{Model of the front-face measurement}\label{sec:frontface}

A closure enters a thermoreflectance experiment through a model of the film response. This section
treats the one-dimensional idealization---plane heating and detection at the front face, no metal
transducer, perfect contact---in which the constitutive law of the film enters through two combinations
of its parameters. It fixes the model, its depth-scaling gauge and the place of the closure coefficient;
the identifiability of the remaining parameters is the subject of a theorem stated in
Sec.~\ref{sec:rank}. The transducer, interface resistances and finite beams of an actual FDTR
experiment are treated in Sec.~\ref{sec:fdtr}.

\subsection{Assumptions}\label{sec:ff-assumptions}

\begin{enumerate}
\item[(F1)] Heat flows in one dimension across a homogeneous film $0<z<e$ on a semi-infinite substrate
$z>e$; the response is linear and the coefficients are constant.
\item[(F2)] The initial state is at rest and there is no interior source. The absorbed flux density
$\wp(t)$ enters at $z=0$, where the temperature rise $\theta(0,t)$ is detected; the front face has no
other losses.
\item[(F3)] In the film, the energy balance $\rho c\,\partial_t\theta+\partial_zq=0$ holds together with the
one-dimensional form of Eq.~\eqref{eq:gk},
\begin{equation}
\tau_R\partial_tq+q=-\lambda\,\partial_z\theta+L^2\partial_z^2q,\qquad L^2=(1+\beta)\,\ell^2 ,
\label{eq:gk1d}
\end{equation}
obtained for $\mathbf q=q(z,t)\,\mathbf e_z$, for which
$\nabla^2\mathbf q+\beta\nabla(\nabla\cdot\mathbf q)=(1+\beta)\,\partial_z^2q\,\mathbf e_z$. The closure
coefficient is $\beta=2$ for the original GK equation and $\beta=1/3$ for the conserving closure;
$L=0$ is Cattaneo's law and $\tau_R=L=0$ Fourier's law.
\item[(F4)] Temperature and normal flux are continuous at $z=e$ (perfect contact, with the
temperature defined by the local energy density on both sides).
\item[(F5)] The substrate obeys Fourier's law with a known effusivity $b_s>0$.
\item[(F6)] The observable is the front-face temperature for all $p=i\omega$, $\omega>0$, unless a finite
set of frequencies is stated.
\end{enumerate}
We write $\theta=T-T_0$, $a=\lambda/\rho c$, $b=(\lambda\rho c)^{1/2}$, $\xi_1=e/\sqrt a$ and
\begin{equation}
\tau_\ell=\frac{L^2}{a}=\frac{(1+\beta)\,\ell^2}{a}.
\label{eq:taul}
\end{equation}
Three-dimensional models of non-Fourier thermoreflectance impose additional slip and temperature-jump
conditions at boundaries~\cite{Beardo2020,HennessyMyers2021}; the one-dimensional problem needs none,
because after elimination the flux in each layer satisfies a second-order equation (below).

\subsection{Film law and the two combinations}\label{sec:combinations}

Laplace transformation of the energy balance from the rest state gives $\rho c\,p\,\theta+q'=0$, hence
$q''=-\rho c\,p\,\theta'$ (primes denote $\partial_z$). Substituting in Eq.~\eqref{eq:gk1d},
$(1+\tau_Rp)\,q=-(\lambda+L^2\rho c\,p)\,\theta'$, that is $q=-\lambda_{\rm eff}(p)\,\theta'$ with
\begin{equation}
\lambda_{\rm eff}(p)=\lambda\,\frac{1+\tau_\ell p}{1+\tau_Rp}.
\label{eq:lameff}
\end{equation}
Then $\theta''=-q'/\lambda_{\rm eff}=\sigma^2\theta$ with $\sigma^2=\rho c\,p/\lambda_{\rm eff}$, and the
flux obeys $[\lambda/(\rho c\,p)+L^2]\,q''=(1+\tau_Rp)\,q$, a second-order equation. The film enters
the solution only through
\begin{equation}
\sigma e=\xi_1\Big[\frac{p(1+\tau_Rp)}{1+\tau_\ell p}\Big]^{1/2},\qquad
\lambda_{\rm eff}\,\sigma=b\,\Big[\frac{p(1+\tau_\ell p)}{1+\tau_Rp}\Big]^{1/2};
\label{eq:combinations}
\end{equation}
the two times exchange their roles between the two combinations.

\subsection{Quadrupole and front-face impedance}\label{sec:quadrupole}

With the flux counted positive toward $z>0$, the state at the front face is related to the state at the
interface by the thermal quadrupole~\PENDING{ref: Maillet \emph{et~al.}, \emph{Thermal Quadrupoles}
(Wiley, 2000); no registered DOI found}~\cite{Krapez2019}
\begin{equation}
\begin{pmatrix}\theta(0)\\ q(0)\end{pmatrix}=
\begin{pmatrix}A&B\\ C&D\end{pmatrix}
\begin{pmatrix}\theta(e)\\ q(e)\end{pmatrix},
\label{eq:quadrupole}
\end{equation}
with $A=D=\cosh\sigma e$, $B=\sinh(\sigma e)/(\lambda_{\rm eff}\sigma)$,
$C=\lambda_{\rm eff}\sigma\sinh\sigma e$ and $AD-BC=1$. These entries are even functions of $\sigma$,
so the branch of the square root is immaterial in the film. The semi-infinite substrate has the
impedance $Z_s=\theta(e)/q(e)=1/(b_s\sqrt p)$ with $\mathrm{Re}\sqrt p>0$, the decaying solution. The
front-face impedance $Z=\theta(0)/q(0)$ is
\begin{equation}
Z(p)=\frac{AZ_s+B}{CZ_s+D}=\frac{1}{\lambda_{\rm eff}\sigma}\,
\frac{\varrho(p)+\tanh\sigma e}{1+\varrho(p)\tanh\sigma e},
\label{eq:impedance}
\end{equation}
where
\begin{equation}
\varrho(p)=\lambda_{\rm eff}\sigma\,Z_s=\frac{b}{b_s}\Big[\frac{1+\tau_\ell p}{1+\tau_Rp}\Big]^{1/2}
\label{eq:rho}
\end{equation}
is a frequency-dependent effusivity ratio; the second form of Eq.~\eqref{eq:impedance} is independent of
the branch chosen for $\sigma$ and remains bounded for thick layers, where $\tanh\sigma e$ saturates.
The front-face temperature is $\hat\theta(0,p)=Z(p)\,\hat\wp(p)$, and in the harmonic regime the
measured amplitude and phase are $|Z(i\omega)|$ and $\arg Z(i\omega)$. The characteristic frequency
$f_c=1/(2\pi\xi_1^2)=a/(2\pi e^2)$ separates the substrate-dominated regime, $f\ll f_c$, from the
film-dominated one.

\subsection{Four combinations and the depth-scaling gauge}\label{sec:gauge}

\begin{proposition}[Dependence and gauge; \status{derived under assumptions}]\label{prop:gauge}
Assume (F1)--(F6); the following statements are exact within this model.
\begin{enumerate}
\item[(i)] $Z(p)$ depends on the film only through $\bm\chi=(\xi_1,b,\tau_R,\tau_\ell)$ and on the
substrate only through $b_s$.
\item[(ii)] For every $\mu>0$ the transformation
\begin{equation}
\begin{aligned}
\mathcal G_\mu:\quad &e\to\mu e,\quad a\to\mu^2a,\quad \lambda\to\mu\lambda,\\
&\rho c\to\rho c/\mu,\quad \ell\to\mu\ell,
\end{aligned}
\label{eq:gaugegroup}
\end{equation}
with $\tau_R$ and $\beta$ fixed, leaves $\bm\chi$, and hence $Z(p)$ for all $p$, invariant.
Infinitesimally, in the variables $(e,a,\lambda,\ell,\tau_R,\beta)$,
\begin{equation}
\big(e\,\partial_e+2a\,\partial_a+\lambda\,\partial_\lambda+\ell\,\partial_\ell\big)\,Z=0 .
\label{eq:euler}
\end{equation}
\item[(iii)] The orbits of $\mathcal G_\mu$ are one-dimensional, and fixing any one of $e$, $a$,
$\lambda$, $\rho c$, or $\ell$ when $\ell>0$, selects a single point on each orbit.
\end{enumerate}
\end{proposition}

\begin{proof}
(i) follows from Eqs.~\eqref{eq:combinations}--\eqref{eq:rho}. (ii) Under $\mathcal G_\mu$,
$\xi_1=e/\sqrt a\to\mu e/(\mu\sqrt a)$, $b=(\lambda\rho c)^{1/2}\to(\mu\lambda\,\rho c/\mu)^{1/2}$ and
$\tau_\ell=(1+\beta)\ell^2/a\to(1+\beta)\mu^2\ell^2/(\mu^2a)$; differentiation at $\mu=1$ gives
Eq.~\eqref{eq:euler}. (iii) $e$, $a$, $\lambda$, $\rho c$ and $\ell$ scale as $\mu$, $\mu^2$, $\mu$,
$\mu^{-1}$ and $\mu$.
\end{proof}

For Fourier films ($\tau_R=\tau_\ell=0$) this is the non-identifiability of thickness, diffusivity and
conductivity shown by Krapez and Rigollet, which holds whatever the frequencies and the noise
level~\cite{KrapezRigollet2017}. Proposition~\ref{prop:gauge} states that neither a relaxation time nor a
nonlocal term removes it, since both enter only through times invariant under depth scaling. The
extension is elementary dimensional analysis; we did not find it stated for non-Fourier laws in the
literature searched. Independent information that fixes a length scale---a measured thickness,
conductivity or heat capacity---selects the point on the orbit; so does an independent velocity $v$
together with the grey kinetic relation $a=v^2\tau_R/3$, conditionally on that relation. Finite beams
with known radii and in-plane properties introduce a lateral length; whether this restores the
identifiability of $(e,a,\lambda)$ is part of the pending analysis of Sec.~\ref{sec:rank}.

\begin{corollary}[Closure coefficient and nonlocal length; \status{derived under assumptions} (F1)--(F6)]\label{cor:beta}
The closure coefficient $\beta$ and the nonlocal length $\ell$ enter $Z$ only through
$\tau_\ell=(1+\beta)\ell^2/a$. Even with $e$ and $a$ known, all pairs $(\beta,\ell)$ on the curve
$(1+\beta)\ell^2=a\tau_\ell$ give the same response. In particular, one-dimensional front-face data do
not distinguish the original ($\beta=2$) from the conserving ($\beta=1/3$) closure when $\ell$ is free,
and the lengths inferred from one fitted $\tau_\ell$ differ by $\ell_{\beta=1/3}/\ell_{\beta=2}=3/2$.
\end{corollary}

Distinguishing the closures therefore requires a multidimensional geometry or independent constitutive
information. Thermodynamics does not fix $\beta$ either: with the entropy density
$s_{\rm eq}-\tau_Rq^2/(2\lambda T_0^2)$ and the entropy flux $q/T+L^2q\,\partial_zq/(\lambda T_0^2)$, the
entropy production to second order in the perturbation is $(q^2+L^2(\partial_zq)^2)/(\lambda T_0^2)\ge0$
for every $L^2\ge0$, within extended irreversible thermodynamics~\cite{JouCasasLebon2010}, and
nonequilibrium-thermodynamic derivations treat the Laplacian and grad-div coefficients as independent
parameters~\cite{Kovacs2024}.

\subsection{Fourier resonance}\label{sec:resonance}

If $\tau_R=\tau_\ell$, then $\lambda_{\rm eff}=\lambda$, $\sigma e=\xi_1\sqrt p$ and
$\lambda_{\rm eff}\sigma=b\sqrt p$, and $Z$ coincides for all $p$ with the response of a Fourier film with
the same $(\xi_1,b)$. This is the Fourier resonance, named by Both \emph{et~al.}~\cite{Both2016} and
analyzed by Kov\'acs~\cite{Kovacs2018}; the equivalence concerns the boundary-value problem
(F1)--(F6) and does not extend to interior sources or arbitrary initial states. Its position in terms of
kinetic times depends on the closure. With the grey coefficients of Sec.~\ref{sec:cattaneo},
$a=v^2\tau_R/3$ and $\ell^2=v^2\tau_c\tau_R/5$, the conserving closure gives $\tau_\ell=4\tau_c/5$ and
\begin{equation}
\frac{\tau_\ell}{\tau_R}=\frac{4\,\tau_N}{5\,(\tau_N+\tau_R)}<\frac45 ,
\label{eq:noresonance}
\end{equation}
so it has no Fourier resonance at positive times, an elementary consequence of the published
closure~\cite{GuoWang2022}. The original closure gives $\tau_\ell=9\tau_c/5$ and a resonance at
$\tau_R/\tau_N=4/5$; with $\tau_c$ replaced by $\tau_N$ it gives $\tau_R/\tau_N=9/5$. Neither ratio
satisfies $\tau_N\ll\tau_R$, which the closure itself assumes, so the resonance line is a
closure-dependent formal marker, not a physical boundary between transport regimes.

\subsection{Limits and checks}\label{sec:ff-limits}

For $\tau_R=\tau_\ell=0$, Eq.~\eqref{eq:impedance} is the classical quadrupole response of a coating
on a substrate, with $\varrho=b/b_s$. For a thick film, $\mathrm{Re}\,\sigma e\to\infty$ and
$Z\to1/(\lambda_{\rm eff}\sigma)$, the impedance of a half-space, with phase
\begin{equation}
\arg Z(i\omega)=-\frac\pi4+\frac12\big[\arctan(\omega\tau_R)-\arctan(\omega\tau_\ell)\big].
\label{eq:halfphase}
\end{equation}
For Cattaneo's law ($\tau_\ell=0$) the phase rises from $-45^\circ$ to $0^\circ$, the closed form of
Camacho de la Rosa \emph{et~al.}~\cite{CamachoDeLaRosa2025}; Eq.~\eqref{eq:halfphase} also agrees with
the one-dimensional limit of the GK model of Hennessy and Myers~\cite{HennessyMyers2021} after
conversion to the time convention $e^{+i\omega t}$. Replacing $p$ by $p(1+\tau_Rp)$ throughout the
Fourier quadrupole gives the correct argument $\sigma e$ for Cattaneo's law but a flux coefficient that
is wrong by the factor $1+\tau_Rp$; the half-space phase then runs from $-45^\circ$ toward $-90^\circ$.
A Cattaneo film propagates disturbances at the finite speed $(a/\tau_R)^{1/2}$~\cite{JosephPreziosi1989};
for $\ell>0$ the term $L^2\partial_z^2q$ makes the principal part parabolic and the propagation speed
unbounded.

\subsection{Identifiability of the four combinations}\label{sec:rank}

Proposition~\ref{prop:gauge} bounds what the measurement can identify: at most $\bm\chi$. Whether the
map $\bm\chi\mapsto Z(\cdot\,;\bm\chi)$ is locally and globally injective---structural identifiability in
the sense of Ref.~\cite{Raue2009}---and on which exceptional set it degenerates, is the following
statement.

\PENDING{Theorem: generic rank four and degenerate set}

% Section V — AlN: film-regularized memory time (pending results)
\section{AlN: film-regularized memory time}\label{sec:aln}

\PENDING{AlN first-principles results: film-regularised memory-time ranges. Planned content:
(i) event geometry, reference operator and validation against an independent solver; infrared control
of the bulk memory moment on the event operator and in the RTA (mesh series), and the status of
hypotheses (H1)--(H5); (ii) the minimum and maximum of $\tmem/\tstat$ in the event-cone class for the
Debye and AlN event sets, with certified witnesses and the constant $M$; (iii) the film-regularized
memory time $\tmem(L)$ of the conserving closure in the force-constant-parametrized rate class at
$30$, $100$ and $300~\mathrm{nm}$, with held-out-mesh checks, the regularization used and its
limitations; every number with its numerical-error estimate and status label}

% Section VI — FDTR detectability (pending results)
\section{Detectability in frequency-domain thermoreflectance}\label{sec:fdtr}

\PENDING{FDTR detectability: phase differences and detectability thresholds of the admissible range of
the film-regularized memory time for the Au/AlN/sapphire scenario ($500~\mathrm{nm}$ AlN,
$10~\mathrm{kHz}$--$200~\mathrm{MHz}$, phase noise $0.01^\circ$--$0.1^\circ$, amplitude noise
$1\,\%$), with the declared nuisance sets (instrument gain, phase and delay; interface resistances;
film conductivity and capacity; nonlocal time and transducer properties; beam radii), convergence of
the Hankel quadrature, frequency grid and Jacobian steps, and an independent one-dimensional check}

% Section VII — Discussion and conclusions (pending)
\section{Discussion and conclusions}\label{sec:conclusions}

\PENDING{conclusions: to be written after Secs.~\ref{sec:aln} and \ref{sec:fdtr} and the rank theorem
are final; restricted to the evidence, with the limitations of the operator classes, of the regularized
memory time and of the one-dimensional measurement model stated explicitly}



\section*{Data availability}
\PENDING{data and code availability statement: repository, archived version and identifiers}

\appendix
% Appendix A — Notation
\section{Notation}\label{app:notation}

Tables~\ref{tab:notation-kinetic} and \ref{tab:notation-film} list the symbols used in the main text.
Vectors in mode space and Cartesian vectors are bold; ${}^{\mathsf T}$ denotes transposition, ${}^+$ the
Moore--Penrose pseudoinverse and $\succeq0$ positive semidefiniteness. The time convention is
$e^{+i\omega t}$, with Laplace variable $p$ and $p=i\omega$ in the harmonic regime; fluxes are densities
per unit area.

\begin{table*}[tbp]
\caption{Kinetic symbols (Secs.~\ref{sec:kinetics} and \ref{sec:identifiability}).}
\label{tab:notation-kinetic}
\footnotesize
\begin{ruledtabular}
\begin{tabular}{p{0.22\textwidth}p{0.58\textwidth}p{0.12\textwidth}}
Symbol & Meaning & Where\\ \hline
$\nu=(\mathbf k,s)$, $n$ & mode index, number of modes & (K1)\\
$N_{\mathbf k}$, $N_1$, $\mathcal V$, $d$ & number of wave vectors, points per direction, crystal volume, spatial dimension & (K1), \ref{sec:ir}\\
$\varepsilon_\nu$, $\mathbf v_\nu$, $f^0_\nu$ & energy, group velocity, Bose occupation & (K1)\\
$w_\nu$, $\mathsf W$ & $[f^0_\nu(1+f^0_\nu)]^{1/2}$, $\mathrm{diag}(w_\nu)$ & (K2)\\
$c_\nu$ & mode heat capacity $\varepsilon_\nu^2w_\nu^2/(k_BT^2)$ & \eqref{eq:harmonic}\\
$\yy$ & entropy coordinates & \eqref{eq:entropycoords}\\
$\Omega$, $\Cop$ & collision operator acting on $\delta f$; the same in entropy coordinates & \eqref{eq:kinetic}\\
$\ee$, $\Hsp$, $n_{\Hsp}$, $\Pi_{\Hsp}$ & energy invariant, its orthogonal complement, $\dim\Hsp=n-1$, orthogonal projector & \eqref{eq:entropy}, (K3)\\
$\jj=\jj^{(i)}$, $q_i$ & current vector, heat-flux density & \eqref{eq:current}\\
$\rho c$, $\Lambda_i$ & volumetric heat capacity, ballistic weight & \eqref{eq:harmonic}\\
$\Phi(t)$, $\kappa(t)$ & current autocorrelation function, memory kernel & \eqref{eq:memorylaw}\\
$K(p)$, $\lambda(p)$, $\lambda$ & response function, dynamic conductivity, dc conductivity & \eqref{eq:dynamicK}\\
$\Gamma_k$, $\mathbf u_k$, $\lambda_k$ & relaxon rates, relaxons, relaxon contributions to $\lambda$ & \eqref{eq:stieltjes}, \eqref{eq:forms}\\
$\xx$, $K_0$, $N$ & $\Cop^+\jj$, $K(0)$, $|\xx|^2$ & \eqref{eq:moments}\\
$\tmem$, $\tstat$ & memory moment $-K'(0)/K(0)$, static time $K_0/|\jj|^2$ & \eqref{eq:moments}\\
$\tau_N$, $\tau_R$, $\tau_c$ & normal, resistive and combined times of the Callaway model & \eqref{eq:callaway}, \eqref{eq:gk}\\
$\mathbf P_i$, $\Pi_{\rm N}$ & crystal momentum, projector onto $\mathrm{span}\{\ee,\mathbf P_x,\mathbf P_y,\mathbf P_z\}$ & \eqref{eq:callaway}\\
$\beta$, $\ell$, $L_{\rm v}$ & closure coefficient, nonlocal length, scale of spatial variation & \eqref{eq:gk}\\
$\mathbf r$, $\tau_\nu$ & mode relaxation rates $\mathrm{diag}\,\Cop$, lifetimes $1/r_\nu$ & Def.~\ref{def:data}\\
$\mathfrak P(\ee)$, $\mathfrak E(Z)$ & positive-semidefinite class, fixed event cone & Def.~\ref{def:classes}\\
$Z=(\bm\zeta_\gamma)_{\gamma\in\mathcal E}$ & event vectors, event set & \eqref{eq:cone}\\
$\mathbf s_\gamma$, $g_\gamma$ & stoichiometric vector, event rate & \eqref{eq:events}\\
$\mathcal F(\mathbf r^\star,K^\star)$ & feasible rates for given static data & \eqref{eq:feasible}\\
$\mathfrak B$, $\yy_I$, $M$ & bases, basic solutions, cone constant & Def.~\ref{def:M}\\
$\bm\psi_I$, $g_I$, $\varpi_I$ & cofactor vector, product of rates, convex weight & \eqref{eq:hull}\\
$\mathsf S_I$, $\bm\phi_I$ & stoichiometric matrix of the events in $I$, hidden invariant & \eqref{eq:subnetwork}\\
$K_{\rm RTA}$, $\varsigma_{\max}$ & RTA dc response, largest number of distinct modes per event & Lemma~\ref{lem:rta}\\
$\Gamma_{\min}$, $g^{\rm ref}$, $F$ & smallest eigenvalue on $\Hsp$, reference rates, prior factor & \eqref{eq:prior}\\
$\alpha$, $r_0$, $J$, $J_0$ & infrared exponent, angular rate coefficient, current density and its long-wavelength limit & (H1)--(H2)\\
$\mathcal M_m$, $k_0$, $k_{\min}$ & RTA moments, infrared radius, cutoff & \eqref{eq:Mm}\\
$\mathsf H$, $\mathbf h_k$, $\mathsf T(\eta)$, $\eta$ & Walsh matrix, its columns, block matrix, family parameter & \eqref{eq:walsh}\\
\end{tabular}
\end{ruledtabular}
\end{table*}

\begin{table*}[tbp]
\caption{Symbols of the front-face model (Sec.~\ref{sec:frontface}).}
\label{tab:notation-film}
\footnotesize
\begin{ruledtabular}
\begin{tabular}{p{0.22\textwidth}p{0.58\textwidth}p{0.12\textwidth}}
Symbol & Meaning & Where\\ \hline
$z$, $e$ & depth, film thickness & (F1)\\
$\theta$, $q$, $\wp$ & temperature rise, flux density (positive toward $z>0$), absorbed flux density & (F2)\\
$\lambda$, $\rho c$, $a$, $b$ & conductivity, volumetric heat capacity, diffusivity, effusivity of the film & \ref{sec:ff-assumptions}\\
$b_s$ & effusivity of the substrate & (F5)\\
$\xi_1$ & $e/\sqrt a$ & \ref{sec:ff-assumptions}\\
$\tau_R$, $\ell$, $\beta$, $L^2$ & relaxation time, nonlocal length, closure coefficient, $(1+\beta)\ell^2$ & \eqref{eq:gk1d}\\
$\tau_\ell$ & nonlocal time $L^2/a$ & \eqref{eq:taul}\\
$\lambda_{\rm eff}$, $\sigma$ & effective conductivity, propagation constant & \eqref{eq:lameff}\\
$A,B,C,D$ & quadrupole entries & \eqref{eq:quadrupole}\\
$Z$, $Z_s$, $\varrho$ & front-face and substrate impedances, frequency-dependent effusivity ratio & \eqref{eq:impedance}, \eqref{eq:rho}\\
$f_c$ & characteristic frequency $1/(2\pi\xi_1^2)$ & \ref{sec:quadrupole}\\
$\bm\chi$ & film combinations $(\xi_1,b,\tau_R,\tau_\ell)$ & Prop.~\ref{prop:gauge}\\
$\mu$, $\mathcal G_\mu$ & gauge parameter, depth-scaling transformation & \eqref{eq:gaugegroup}\\
\end{tabular}
\end{ruledtabular}
\end{table*}


\bibliography{refs}

\end{document}
